\documentclass[11pt]{article}

\usepackage[T1]{fontenc}
\usepackage[utf8]{inputenc}
\usepackage{lmodern}
\usepackage{microtype}
\usepackage[a4paper,margin=28mm]{geometry}
\usepackage{amsmath,amssymb,amsthm,mathtools}
\usepackage{booktabs,array}
\usepackage[dvipsnames]{xcolor}
\usepackage{enumitem}
\usepackage{hyperref}
\hypersetup{
 colorlinks=true,
 linkcolor=NavyBlue,
 urlcolor=NavyBlue,
 citecolor=NavyBlue,
 pdftitle={Spectral Norms of Independent-Entry Matrices with Regular Moment Growth},
 pdfauthor={Diar Heidary},
 pdfkeywords={random matrices, spectral norm, regular moment growth, MI-32, positive polynomials}
}
\setlist{itemsep=0.25em,topsep=0.35em}
\setlength{\emergencystretch}{2em}
\allowdisplaybreaks

\newtheorem{theorem}{Theorem}[section]
\newtheorem{lemma}[theorem]{Lemma}
\newtheorem{proposition}[theorem]{Proposition}
\newtheorem{corollary}[theorem]{Corollary}
\theoremstyle{definition}
\newtheorem{definition}[theorem]{Definition}
\newtheorem{example}[theorem]{Example}
\theoremstyle{remark}
\newtheorem{remark}[theorem]{Remark}
\newcommand{\R}{\mathbb R}
\newcommand{\C}{\mathbb C}
\newcommand{\E}{\mathbb E}
\newcommand{\Prob}{\mathbb P}
\newcommand{\norm}[1]{\lVert#1\rVert}
\newcommand{\ip}[2]{\langle#1,#2\rangle}
\newcommand{\tr}{\operatorname{tr}}
\newcommand{\diag}{\operatorname{diag}}

\title{Spectral Norms of Independent-Entry Matrices\\with Regular Moment Growth}
\author{Diar Heidary}
\date{Research draft --- 14 September 2026}

\begin{document}
\maketitle

\begingroup
\setlength{\fboxrule}{1pt}
\setlength{\fboxsep}{8pt}
\noindent
\fcolorbox{NavyBlue!75!black}{NavyBlue!3}{%
 \begin{minipage}{\dimexpr\linewidth-2\fboxsep-2\fboxrule\relax}
 \small\raggedright
 \textbf{\color{NavyBlue}AI-assisted draft; the upper bound is
 machine-checked.}
 The upper comparison of Theorem~\ref{mi32-main-theorem} has been
 formalized in Lean~4 and checked by the Lean kernel. The formal target
 carries no proof gap, and its transitive axiom audit reports only
 \texttt{propext}, \texttt{Classical.choice} and \texttt{Quot.sound};
 Section~\ref{mi32-formalization} records the exact statement, the
 development, and the two places where the formal proof deliberately
 differs from the informal one. The lower comparison in
 Section~\ref{mi32-lower-section} is a known result of
 Lata\l a--\'Swi\k atkowski, reproduced here for completeness, and is
 \emph{not} formalized. AI tools assisted the derivation, the
 formalization and the exposition under the author's direction. There
 has been no external peer review. The upper comparison is registered
 with the Palomar registry as entry
 	exttt{PALOMAR-2026-09-14-000006}, whose mechanical verification
 replayed the development through the Comparator and NanoDa kernels;
 registration is not acceptance, endorsement, or peer review.
 \end{minipage}%
}
\endgroup
\medskip

\begin{abstract}
We prove the upper comparison in problem MI-32, that is, the upper half
of Conjecture~4.3 of Lata\l a--\'Swi\k atkowski: for real matrices with
independent mean-zero entries whose moments satisfy
$\norm{X_{ij}}_{L_{2p}}\le\alpha\norm{X_{ij}}_{L_p}$, the expected
operator norm is at most $C_\alpha$ times the sum of the row and column
variance scales and the shared-index deletion parameter built from weak
bilinear moments at order $\log(k+1)$, including the exponent
$\log2<1$ at $k=1$.

The local estimate comes from representing each entry as an independent
sign times its magnitude. Sign parity supplies orthogonal sectors while
every magnitude power stays inside the polynomial coefficients, so
revisiting an entry does not lose its moment information. Creation
localizes a few original rows and annihilation a few original columns,
and a net on whichever axis is small costs an absolute factor at the
relevant moment order. The remaining ingredient is a moment comparison
on the cone of polynomials with nonnegative original monomial
coefficients, which contains every vacuum walk polynomial and is
preserved by the parity projections: in degree $d$ such a polynomial
obeys an $L_2$-to-$L_4$ bound with factor $(\sqrt3\,\alpha^2)^d$.

The argument is self-contained. In place of the strong/weak moment
comparison for Banach-space-valued sums we prove the
independent-coordinate statement the matrix argument actually needs, by
an independent copy and a multinomial expansion; in place of the
variance-envelope matrix estimate used in the published reduction of the
deletion parameter we raise the screening threshold from $b^2$ to $b^3$
and estimate the far remainder by an elementary centered Gram second
moment. All constants are explicit, and all of them are constants that
a proof assistant has checked. We also reproduce the established lower
comparison, with attribution, and record precisely what the formal
development does and does not certify.
\end{abstract}

\clearpage
\begingroup
\small
\tableofcontents
\endgroup
\clearpage

% MI32 BODY BEGIN


\section{Introduction and the exact mean-norm statement}
\label{mi32-introduction}

Let $X=(X_{ij})_{i,j\le n}$ be a real matrix with independent mean-zero
entries. We assume that every absolute moment is finite and that one
constant $\alpha\ge1$ satisfies
\begin{equation}\label{mi32-regularity}
 \norm{X_{ij}}_{L_{2p}}\le\alpha\norm{X_{ij}}_{L_p},
 \qquad p\ge1,\quad i,j\in[n].
\end{equation}
Here $[n]=\{1,\ldots,n\}$ and
$\norm{Z}_{L_p}=(\E|Z|^p)^{1/p}$ for every $p>0$. At $p<1$ this
notation denotes a moment functional, not a norm. All logarithms below
are natural. Matrix norms without a subscript are operator norms for
the Euclidean norms on domain and range. For a scalar random variable
we also use the abbreviation $\norm{Z}_p=\norm{Z}_{L_p}$; for a
deterministic vector, the subscript $2$ denotes its Euclidean norm.

Define the deterministic variance scales
\begin{equation}\label{mi32-variances}
 \begin{aligned}
 \sigma_r(X)&=\max_i\Big(\sum_j\E X_{ij}^2\Big)^{1/2},&
 \sigma_c(X)&=\max_j\Big(\sum_i\E X_{ij}^2\Big)^{1/2},\\
 M(X)&=\sigma_r(X)+\sigma_c(X),&
 \sigma(X)&=\max\{\sigma_r(X),\sigma_c(X)\}.
 \end{aligned}
\end{equation}
For a deterministic set $I\subset[n]$ and $p>0$, put
\begin{equation}\label{mi32-weak}
 R_{X,I}(p)=\sup_{\norm{s}_2,\norm{t}_2\le1}
     \norm{\sum_{i,j\notin I}X_{ij}s_it_j}_{L_p},
 \qquad R_X(p)=R_{X,\varnothing}(p).
\end{equation}
The vectors in this supremum are deterministic. The deletion parameter is
\begin{equation}\label{mi32-D}
 D(X)=\max_{1\le k\le n}\ \min_{I\subset[n],\ |I|\le k}
                                  R_{X,I}(\log(k+1)).
\end{equation}
One and the same original set $I$ deletes both rows and columns. It is
chosen for the law, before any realization is observed. The minimum is
outside the vector supremum and is attained, since the collection of
sets is finite. Deleting all indices gives zero, so the $k=n$ term is
zero. The $k=1$ term retains its exact exponent $\log2<1$.

\begin{theorem}[Independent entries with regular moment growth]
\label{mi32-main-theorem}
For every $\alpha\ge1$ there are constants $c_\alpha,C_\alpha>0$
such that every $n\ge1$ and every matrix satisfying
\eqref{mi32-regularity} obey
\begin{equation}\label{mi32-main-bound}
 c_\alpha\{M(X)+D(X)\}\le\E\norm{X}
                       \le C_\alpha\{M(X)+D(X)\}.
\end{equation}
The constants depend only on $\alpha$, not on the dimension, the
individual entry distributions, or the deterministic deletion sets.
\end{theorem}

Only the upper comparison in \eqref{mi32-main-bound} is proved here. It
is the new content of this paper, and it is the half that has been
machine-checked; an explicit admissible value of $C_\alpha$ is recorded
in \eqref{mi32-final-constant}. The lower comparison is a known result,
reproduced with attribution in Section~\ref{mi32-lower-section}.
Section~\ref{mi32-literature} locates both halves in the literature and
says precisely what the present argument adds.

\subsection{The proof mechanism}

For symmetric entries, write $Y_e=\varepsilon_e Z_e$ in law, where
$Z_e=|Y_e|$ and the original signs are independent of all magnitudes.
The signs give an orthogonal parity basis. Magnitude powers stay inside
its coefficients, so visiting an entry repeatedly does not lose its
moment information. In the row-to-column direction, parity creation
has exact sectors involving few original rows, whereas parity
annihilation has different exact sectors involving few original
columns. A net on this small axis has constant cost at the relevant
moment order.

The remaining estimate is a polynomial moment comparison on the cone
of nonnegative original monomial coefficients. That cone contains
every vacuum walk polynomial and is preserved by the parity projections
and every subsequent original multiplication. It is not the entire
linear span of these polynomials. This distinction permits the argument
to retain rare-amplitude effects without bounding an unrestricted
orthogonal-polynomial multiplication operator.

Every analytic input is proved below from \eqref{mi32-regularity}: the
Boolean fourth-moment inequality
(Lemma~\ref{mi32-hilbert-boolean}), the positive-cone comparison
(Lemma~\ref{mi32-cone-lemma}), the scalar and vector moment comparisons
of Section~\ref{mi32-comparison-section}, and the deletion reduction of
Section~\ref{mi32-reduction-section}. See
Remark~\ref{mi32-no-external-inputs}.

\subsection{Relation to the literature}
\label{mi32-literature}

Estimate \eqref{mi32-main-bound} is the content of
\cite[Conjecture~4.3]{LatalaSwiatkowski}, recorded as problem MI-32 in
\cite{MI32}. Its lower half is
\cite[Theorem~4.1]{LatalaSwiatkowski}, which we reproduce in
Section~\ref{mi32-lower-section}; the upper half is what is proved
here. The conjecture refines a long line of work on the norm of a
matrix with independent but not identically distributed entries:
Seginer's comparison with the largest row and column
\cite{Seginer}, Lata\l a's variance-based bound \cite{Latala2005}, the
sharp Gaussian and sub-Gaussian estimates of Bandeira--van~Handel
\cite{BandeiraVanHandel}, the logarithmic-factor results of
Riemer--Sch\"utt \cite{RiemerSchutt}, and the dimension-free structure
theory of Lata\l a--van~Handel--Youssef \cite{LatalaVanHandelYoussef};
see \cite{VanHandelSurvey} for a survey. For weighted Rademacher
matrices the conjecture is known up to a triple logarithm
\cite{Latala2024}, and for general independent entries up to an
iterated logarithm \cite{Meller}; the present estimate removes the
remaining gap under the regularity hypothesis
\eqref{mi32-regularity}.

The relation to \cite{LatalaSwiatkowski} is worth stating precisely,
because it isolates exactly what is new.
\cite[Remark~4.5]{LatalaSwiatkowski} shows that, to obtain the upper
half of their conjecture, it suffices to prove the \emph{undeleted}
local estimate
\begin{equation}\label{mi32-lsw-sufficient}
 \E\norm{X}\le C_\alpha\bigl\{M(X)+R_X(\log(n+1))\bigr\}.
\end{equation}
Corollary~\ref{mi32-local-mean} below is exactly
\eqref{mi32-lsw-sufficient} for entries symmetric in law, and it is the
new local estimate of this paper. The deduction of the conjecture from
\eqref{mi32-lsw-sufficient} sketched in
\cite[Remark~4.5]{LatalaSwiatkowski} passes through a variance-envelope
bound taken from \cite[Theorem~4.4]{LatalaVanHandelYoussef}.
Section~\ref{mi32-reduction-section} gives a different deduction that
avoids that input; the reason and the cost are explained there and in
Remark~\ref{mi32-cube-remark}.

\subsection{Organization}

Sections~\ref{mi32-polynomial-section} and
\ref{mi32-comparison-section} prove the polynomial, scalar and vector
moment estimates used later. Section~\ref{mi32-representation-section}
sets up the parity--amplitude representation and
Section~\ref{mi32-counts-section} the two count decompositions.
Sections~\ref{mi32-thin-section} and \ref{mi32-vacuum-section} prove the
local moment theorem and its consequence
\eqref{mi32-lsw-sufficient} at the exact order $\log(n+1)$.
Section~\ref{mi32-reduction-section} replaces the undeleted weak moment
by the deletion parameter and removes the symmetry assumption, which
completes the upper bound. Section~\ref{mi32-lower-section} reproduces
the known lower comparison, Section~\ref{mi32-scope} delimits what is
claimed, and Section~\ref{mi32-formalization} describes the machine
verification.
\section{Polynomial moment estimates from the entry law}
\label{mi32-polynomial-section}

Every analytic ingredient of the upper bound is proved in this section
and the next, from the entry hypothesis \eqref{mi32-regularity} alone.
Two features of this are worth stating at the outset, because earlier
treatments of the same route depend on external matrix theorems.

\begin{remark}[No external analytic input]
\label{mi32-no-external-inputs}
The upper bound uses no strong/weak moment comparison for
Banach-space-valued sums, and no variance-envelope matrix estimate. The
vector comparison the argument needs is the
independent-coordinate statement of
Theorem~\ref{mi32-hilbert-comparison}, which is proved here by an
independent copy and a multinomial expansion. The far remainder of the
deletion argument is estimated in
Section~\ref{mi32-far-remainder} by an elementary centered Gram second
moment. Consequently the only results quoted from the literature are the
lower comparison of Section~\ref{mi32-lower-section}, which is not part
of the estimate proved here, and the classical Boolean fourth-moment
inequality of Lemma~\ref{mi32-hilbert-boolean}, for which we give a
complete proof.
\end{remark}

\subsection{Hilbert-valued Boolean fourth moments}

On independent signs $\varepsilon_1,\ldots,\varepsilon_N$, write
$\varepsilon_S=\prod_{e\in S}\varepsilon_e$ and let $N_\rho$ be the
noise operator defined on Fourier characters by
$N_\rho\varepsilon_S=\rho^{|S|}\varepsilon_S$. Only the
$L_2$-to-$L_4$ case of Boolean hypercontractivity is needed, and only
for Hilbert-space-valued coefficients. We give its short proof; the
general Boolean theorem is in \cite[Ch.~9--10]{ODonnellBoolean}.

\begin{lemma}[Hilbert-valued Boolean polynomials]
\label{mi32-hilbert-boolean}
If $F=\sum_{|S|\le d}\varepsilon_S v_S$ takes values in a real
Euclidean space, then
\begin{equation}\label{mi32-hilbert-boolean-bound}
 \norm{F}_{L_4(\ell_2)}\le3^{d/2}\norm{F}_{L_2(\ell_2)}.
\end{equation}
\end{lemma}
\begin{proof}
For fixed real Hilbert vectors $a,b$,
\[
 \E_\varepsilon\norm{a+\varepsilon b}_2^4
 =(\norm{a}_2^2+\norm{b}_2^2)^2+4\ip{a}{b}^2
 \le\norm{a}_2^4+6\norm{a}_2^2\norm{b}_2^2+\norm{b}_2^4.
\]
If $A,B$ depend on other signs, integrate this inequality and use
Cauchy--Schwarz in the mixed term. Taking a square root gives
\begin{equation}\label{mi32-one-sign-fourth}
 \norm{A+\varepsilon B}_{L_4(\ell_2)}^2
 \le\norm{A}_{L_4(\ell_2)}^2+3\norm{B}_{L_4(\ell_2)}^2.
\end{equation}
Set $\rho=1/\sqrt3$ and induct on the number of signs. Write
$F=A+\varepsilon_N B$, where $A,B$ depend on the first $N-1$ signs.
Denoting by $N_\rho'$ the noise operator on those signs, the induction
hypothesis and \eqref{mi32-one-sign-fourth} yield
\[
 \norm{N_\rho F}_{L_4(\ell_2)}^2
 \le\norm{N_\rho'A}_{L_4(\ell_2)}^2
        +3\rho^2\norm{N_\rho'B}_{L_4(\ell_2)}^2
 \le\norm{A}_{L_2(\ell_2)}^2+\norm{B}_{L_2(\ell_2)}^2
 =\norm{F}_{L_2(\ell_2)}^2.
\]
The base case of zero signs is equality. Finally set
$G=\sum_{|S|\le d}\rho^{-|S|}\varepsilon_Sv_S$, so $N_\rho G=F$.
Hilbert Fourier orthogonality gives
\[
 \norm{F}_{L_4(\ell_2)}
 \le\norm{G}_{L_2(\ell_2)}
 =\Big(\sum_{|S|\le d}\rho^{-2|S|}\norm{v_S}_2^2\Big)^{1/2}
 \le\rho^{-d}\norm{F}_{L_2(\ell_2)}.
\]
\end{proof}

\subsection{Product moments of the magnitudes}

Fix independent symmetric variables $W=(W_e)_{e\le N}$ satisfying
\eqref{mi32-regularity}, and put $Z_e=|W_e|$. For an exponent vector
$\nu\in\mathbb N_0^N$, write $Z^\nu=\prod_e Z_e^{\nu_e}$ and
$|\nu|=\sum_e\nu_e$. Polynomial degree below always means this
ordinary total degree, including all repeated powers.

\begin{lemma}[A product-moment bound]
\label{mi32-product-moment-lemma}
For $m_e(u)=\E Z_e^u$, $u\in\mathbb N_0$, and $m_e(0)=1$,
\begin{equation}\label{mi32-product-moment}
 m_e(u+v)\le\alpha^{u+v}m_e(u)m_e(v)
          \le\alpha^{2(u+v)}m_e(u)m_e(v),
 \qquad u,v\in\mathbb N_0.
\end{equation}
Consequently, every scalar polynomial
$G(Z)=\sum_\nu a_\nu Z^\nu$ with $a_\nu\ge0$ and degree at most
$D$ satisfies
\begin{equation}\label{mi32-magnitude-polynomial}
 \E G(Z)^2\le\alpha^{4D}(\E G(Z))^2.
\end{equation}
\end{lemma}
\begin{proof}
For positive integers $u,v$, Cauchy--Schwarz and moment doubling give
\[
 \begin{aligned}
 m_e(u+v)&\le\sqrt{m_e(2u)m_e(2v)},\\
 m_e(2u)&=\norm{W_e}_{L_{2u}}^{2u}
        \le\alpha^{2u}\norm{W_e}_{L_u}^{2u}
         =\alpha^{2u}m_e(u)^2.
 \end{aligned}
\]
Apply the second formula also with $v$ to prove the first inequality in
\eqref{mi32-product-moment}; the second holds because $\alpha\ge1$.
If either exponent is zero, the assertion follows from $\alpha\ge1$;
this does not apply moment doubling at zero. Identically zero variables
cause no exception.

For two monomials of $G$, independence and the weaker form of
\eqref{mi32-product-moment} give
\[
 \E Z^{\nu+\mu}
 =\prod_e m_e(\nu_e+\mu_e)
 \le\alpha^{2(|\nu|+|\mu|)}
             \prod_e m_e(\nu_e)\prod_e m_e(\mu_e).
\]
Multiply by $a_\nu a_\mu\ge0$, sum over $\nu,\mu$, and use
$|\nu|+|\mu|\le2D$ to obtain \eqref{mi32-magnitude-polynomial}.
\end{proof}

\begin{lemma}[Positive-coefficient polynomial comparison]
\label{mi32-cone-lemma}
Let $F(W)$ be a polynomial taking values in a fixed finite-dimensional
real Euclidean space. Suppose every coordinate has nonnegative
coefficients in the ordinary monomial basis of the original $W_e$,
and $\deg F\le d$. With
\begin{equation}\label{mi32-H-constant}
 H_\alpha=\sqrt3\,\alpha^2,
\end{equation}
one has
\begin{align}
 \norm{F}_{L_4(\ell_2)}&\le H_\alpha^d\norm{F}_{L_2(\ell_2)},
                  \label{mi32-cone-four}\\
 \norm{F}_{L_{2p/(p-2)}(\ell_2)}
   &\le H_\alpha^{4d/p}\norm{F}_{L_2(\ell_2)},\qquad p\ge4.
                  \label{mi32-cone-holder}
\end{align}
\end{lemma}
\begin{proof}
Represent $W_e=\varepsilon_eZ_e$ in law with independent signs and
independent magnitudes. Group ordinary monomials by their exponent
parity:
\[
 F(\varepsilon Z)=\sum_S\varepsilon_SF_S(Z).
\]
Every coordinate of each $F_S$ still has nonnegative coefficients.
Its monomials have exponent parity $S$, degree at most $d$, and hence
$|S|\le d$. Conditional on $Z$, Lemma~\ref{mi32-hilbert-boolean}
gives
\begin{equation}\label{mi32-conditional-fourth}
 \E_\varepsilon\norm{F(\varepsilon Z)}_2^4
 \le3^{2d}\Big(\sum_S\norm{F_S(Z)}_2^2\Big)^2.
\end{equation}
The scalar polynomial $G(Z)=\sum_S\norm{F_S(Z)}_2^2$ has
nonnegative coefficients and degree at most $2d$. Conditional sign
orthogonality gives $\E G=\E\norm{F}_2^2$. Thus
\eqref{mi32-magnitude-polynomial}, with $D=2d$, proves
\[
 \E\norm{F}_2^4\le3^{2d}\alpha^{8d}(\E\norm{F}_2^2)^2
                =H_\alpha^{4d}(\E\norm{F}_2^2)^2.
\]
Its fourth root is \eqref{mi32-cone-four}. For
$u=2p/(p-2)\in[2,4]$, the interpolation identity
$1/u=(1-4/p)/2+(4/p)/4$ proves \eqref{mi32-cone-holder}.
Degree zero gives equality in both comparisons.
\end{proof}

\begin{remark}[A sharper cone constant, not used]
\label{mi32-sharp-cone-remark}
The first inequality of \eqref{mi32-product-moment} is stronger than
the one used, and yields $\E G^2\le\alpha^{2D}(\E G)^2$ in place of
\eqref{mi32-magnitude-polynomial}, hence $H_\alpha=\sqrt3\,\alpha$.
We deliberately carry the conservative value $\sqrt3\,\alpha^2$,
because it is the one the formal development certifies
(Section~\ref{mi32-formalization}); every constant displayed in this
manuscript is therefore a machine-checked constant. Using the sharp
value replaces one factor $\alpha^2$ by $\alpha$ at each occurrence of
$H_\alpha$ below, and changes nothing else.
\end{remark}

The hypothesis of Lemma~\ref{mi32-cone-lemma} concerns coefficients,
not pointwise positivity of $F(W)$. A coordinate such as $W_1+W_2$ is
allowed although its value may be negative. Nor does the lemma apply to
every polynomial merely because the law is symmetric: arbitrary negative
monomial coefficients may create cancellations in the $L_2$ norm. We
will only project or multiply polynomials in ways that preserve the
stated coefficient cone.

\subsection{Scalar linear forms double at every even order}

The following corollary is the only consequence of
Lemma~\ref{mi32-cone-lemma} needed for scalar sums. It is what replaces
the scalar moment-growth estimate that an external strong/weak
comparison would otherwise supply.

\begin{corollary}[Even-order doubling for linear forms]
\label{mi32-scalar-doubling}
Let $W=(W_e)$ be independent symmetric variables satisfying
\eqref{mi32-regularity}, let $a_e$ be arbitrary real coefficients, and
put $S=\sum_ea_eW_e$. For every integer $q\ge1$,
\begin{equation}\label{mi32-scalar-doubling-bound}
 \norm{S}_{L_{4q}}\le H_\alpha\norm{S}_{L_{2q}}.
\end{equation}
\end{corollary}
\begin{proof}
Represent $W_e=\varepsilon_eZ_e$. Replacing $a_e$ by $|a_e|$ changes
neither side: the signs $\mathrm{sgn}(a_e)\varepsilon_e$ have the same
joint law as the $\varepsilon_e$, and both moments depend only on that
law. So assume $a_e\ge0$. Then $S^q$ is a scalar polynomial in the
original variables of degree $q$ whose ordinary coefficients are
nonnegative, being multinomial coefficients times products of the
$a_e$. Lemma~\ref{mi32-cone-lemma}, applied to the one-dimensional
Euclidean space with $d=q$, gives
\[
 \E S^{4q}=\E|S^q|^4\le H_\alpha^{4q}\bigl(\E|S^q|^2\bigr)^2
               =H_\alpha^{4q}\bigl(\E S^{2q}\bigr)^2 .
\]
Taking $4q$-th roots gives \eqref{mi32-scalar-doubling-bound}.
\end{proof}

\section{A strong/weak comparison for independent coordinates}
\label{mi32-comparison-section}

The thin-matrix estimate of Section~\ref{mi32-thin-section} needs to
compare the $L_{4q}$ moment of the Euclidean norm of a random vector
with its variance and its weak moments. For a vector with
\emph{independent} coordinates this can be proved directly, and that is
all the matrix argument uses. We do so here. No comparison for general
Banach norms or general index sets is required, and none is assumed.

\subsection{An even-moment comparison for independent symmetric families}

\begin{lemma}[Coordinatewise even moments tensorize]
\label{mi32-even-comparison}
Let $(A_j)_{j\in J}$ and $(B_j)_{j\in J}$ be finite families of
symmetric real random variables, each family independent, with moments
through order $2q$ finite. The two families may live on different
probability spaces, or on one space with arbitrary dependence between
them. If $c\ge0$ and
\[
 \E A_j^{2r}\le c^{2r}\,\E B_j^{2r}
 \qquad(j\in J,\ 1\le r\le q),
\]
then
\begin{equation}\label{mi32-even-comparison-bound}
 \E\Bigl(\sum_jA_j\Bigr)^{2q}
 \le c^{2q}\,\E\Bigl(\sum_jB_j\Bigr)^{2q}.
\end{equation}
\end{lemma}
\begin{proof}
Expand the $2q$-th power by the multinomial formula. Independence
factors each monomial expectation as $\prod_j\E A_j^{k_j}$ with
$\sum_jk_j=2q$; every such monomial is integrable, by independence and
finiteness of the individual powers. Symmetry makes $\E A_j^{k_j}=0$
whenever $k_j$ is odd, so only terms with all $k_j$ even survive. For
those terms every factor is nonnegative and the hypothesis applies
coordinatewise; the multinomial coefficients are nonnegative integers,
and the powers of $c$ multiply to $c^{\sum_jk_j}=c^{2q}$. Summing the
surviving terms gives \eqref{mi32-even-comparison-bound}. Exponent-zero
factors equal one by normalization.
\end{proof}

Note what this argument does \emph{not} do: it does not compare
expansions of absolute values. The odd monomials vanish exactly, and
the remaining ones are compared term by term.

\subsection{The comparison}

\begin{theorem}[Independent-coordinate strong/weak comparison]
\label{mi32-hilbert-comparison}
Let $(V_j)_{j\in J}$ be a finite family of independent real random
variables, each symmetric in law, with finite moments of order $4q$ for
an integer $q\ge1$. Suppose $\beta\ge0$ satisfies
\begin{equation}\label{mi32-coordinate-regularity}
 \norm{V_j}_{L_{4r}}\le\beta\norm{V_j}_{L_{2r}}
 \qquad(j\in J,\ 1\le r\le q).
\end{equation}
Put
\[
 R=\Bigl(\sum_jV_j^2\Bigr)^{1/2},\qquad
 \sigma^2=\sum_j\E V_j^2,\qquad
 W_s=\sup_{\sum_jt_j^2\le1}
        \Bigl\|\sum_jt_jV_j\Bigr\|_{L_s}.
\]
Then
\begin{equation}\label{mi32-hilbert-comparison-bound}
 \norm{R}_{L_{4q}}\le\sigma+2\beta^2W_{2q}.
\end{equation}
\end{theorem}
\begin{proof}
All the displayed quantities are finite: $R\le\sum_j|V_j|$ pointwise,
Minkowski controls the latter at every order used, and every unit-ball
coefficient satisfies $|t_j|\le1$, so
$W_s\le\sum_j\norm{V_j}_{L_s}<\infty$. The supremum defining $W_s$ is a
supremum of deterministic numbers; no measurable selection is involved.
For $J=\varnothing$ every quantity vanishes.

Work on a product space carrying an independent copy $(V_j')$ of the
whole family, and set
\[
 A_j=V_j^2-(V_j')^2,\qquad B_j=V_jV_j'.
\]
The pairs $(V_j,V_j')$ are independent in $j$, so $(A_j)$ and $(B_j)$
are independent families. Each $A_j$ is symmetric by exchanging the two
copies, and each $B_j$ is symmetric because the law of $V_j$ is
symmetric and the copy is independent. These are statements about laws;
no sign-flipping map on the underlying space is needed.

For $1\le r\le q$, Minkowski, \eqref{mi32-coordinate-regularity} and
exact factorization of the two independent copies give
\begin{equation}\label{mi32-AB-comparison}
 \norm{A_j}_{L_{2r}}\le2\norm{V_j^2}_{L_{2r}}
   =2\norm{V_j}_{L_{4r}}^2
   \le2\beta^2\norm{V_j}_{L_{2r}}^2
   =2\beta^2\norm{B_j}_{L_{2r}}.
\end{equation}
Raising \eqref{mi32-AB-comparison} to the power $2r$ and applying
Lemma~\ref{mi32-even-comparison} with $c=2\beta^2$ yields
\begin{equation}\label{mi32-AB-sums}
 \Bigl\|\sum_jA_j\Bigr\|_{L_{2q}}
 \le2\beta^2\Bigl\|\sum_jB_j\Bigr\|_{L_{2q}}.
\end{equation}

Put $S=\sum_jV_j^2$ and $S'=\sum_j(V_j')^2$, so $\E S'=\sigma^2$.
Jensen in the copy coordinate, applied to $z\mapsto|z|^{2q}$ and then
integrated in the first coordinate, gives
\begin{equation}\label{mi32-jensen-copy}
 \norm{S-\sigma^2}_{L_{2q}}\le\norm{S-S'}_{L_{2q}}
 =\Bigl\|\sum_jA_j\Bigr\|_{L_{2q}} .
\end{equation}
For a fixed deterministic vector $v'$, homogeneity of the definition of
$W_{2q}$ gives $\norm{\sum_jv_j'V_j}_{L_{2q}}\le W_{2q}\norm{v'}_2$;
this is trivial at $v'=0$ and otherwise follows by normalizing.
Raising to the power $2q$, integrating over the independent copy and
taking the root,
\begin{equation}\label{mi32-B-weak}
 \Bigl\|\sum_jB_j\Bigr\|_{L_{2q}}\le W_{2q}\norm{R}_{L_{2q}}.
\end{equation}
This is the integration of a deterministic scalar bound; the inner law
is always the original law of $V$, and the coefficient vector is the
independent copy.

Let $x=\norm{R}_{L_{4q}}$ and $y=\norm{R}_{L_{2q}}$, both finite and
nonnegative. Since $\norm{S}_{L_{2q}}=x^2$, combining Minkowski with
\eqref{mi32-jensen-copy}, \eqref{mi32-AB-sums} and
\eqref{mi32-B-weak}, and then $y\le x$, gives
\[
 x^2\le\sigma^2+2\beta^2W_{2q}\,y\le\sigma^2+2\beta^2W_{2q}\,x .
\]
For nonnegative $x,\sigma,a$ the inequality $x^2\le\sigma^2+ax$ forces
$x\le\sigma+a$: otherwise $x>\sigma+a\ge0$ would give
$x^2>(\sigma+a)x\ge\sigma^2+ax$. Taking $a=2\beta^2W_{2q}$ proves
\eqref{mi32-hilbert-comparison-bound}.
\end{proof}

The hypotheses of Theorem~\ref{mi32-hilbert-comparison} are finite-order
and concern only the individual coordinates. They contain neither the
conclusion nor any vector moment estimate, and there are no dimension
factors in \eqref{mi32-hilbert-comparison-bound}.

\subsection{Doubling at every real order above \texorpdfstring{$\log2$}{log 2}}

The deletion argument of Section~\ref{mi32-reduction-section} compares
weak moments at the subunit order $\log2$ with weak moments at large
scheduled orders. The next proposition supplies that comparison from
Corollary~\ref{mi32-scalar-doubling} alone, with an explicit constant
and no restriction to orders above one.

\begin{proposition}[Uniform real-order doubling]
\label{mi32-real-doubling}
Let $S=\sum_ea_eW_e$ be as in Corollary~\ref{mi32-scalar-doubling} and
put
\begin{equation}\label{mi32-J-constant}
 J_\alpha=H_\alpha^7=(\sqrt3\,\alpha^2)^7 .
\end{equation}
Then
\begin{equation}\label{mi32-real-doubling-bound}
 \norm{S}_{L_{2p}}\le J_\alpha\norm{S}_{L_p},
 \qquad p\ge\tfrac12 .
\end{equation}
In particular \eqref{mi32-real-doubling-bound} holds for every
$p\ge\log2$.
\end{proposition}
\begin{proof}
Write $\beta=H_\alpha\ge1$. Corollary~\ref{mi32-scalar-doubling} at
$q=1$ gives $\norm{S}_{L_4}\le\beta\norm{S}_{L_2}$, that is,
$\E|S|^4\le\beta^4(\E|S|^2)^2$. H\"older with the interpolation
$2=\tfrac47\cdot\tfrac12+\tfrac37\cdot4$ gives
\[
 \E|S|^2\le\bigl(\E|S|^{1/2}\bigr)^{4/7}
              \bigl(\E|S|^4\bigr)^{3/7},
\]
hence $(\E|S|^2)^7\le(\E|S|^{1/2})^4(\E|S|^4)^3
\le\beta^{12}(\E|S|^{1/2})^4(\E|S|^2)^6$, so
$\E|S|^2\le\beta^{12}(\E|S|^{1/2})^4$ and therefore
\begin{equation}\label{mi32-two-to-half}
 \norm{S}_{L_2}\le\beta^6\norm{S}_{L_{1/2}} .
\end{equation}
All these integrals are finite, and the argument is unaffected if one
of them vanishes.

Let $\tfrac12\le p\le2$. Monotonicity of $p\mapsto\norm{S}_{L_p}$ on a
probability space, Corollary~\ref{mi32-scalar-doubling} at $q=1$ and
\eqref{mi32-two-to-half} give
\[
 \norm{S}_{L_{2p}}\le\norm{S}_{L_4}
   \le\beta\norm{S}_{L_2}
   \le\beta^7\norm{S}_{L_{1/2}}
   \le\beta^7\norm{S}_{L_p} .
\]
Let $p>2$ and put $r=\lfloor p/2\rfloor\ge1$, so that
$2r\le p\le4r$ and $2p\le8r$. Two applications of
Corollary~\ref{mi32-scalar-doubling}, at $2r$ and at $r$, give
\[
 \norm{S}_{L_{2p}}\le\norm{S}_{L_{8r}}
   \le\beta\norm{S}_{L_{4r}}
   \le\beta^2\norm{S}_{L_{2r}}
   \le\beta^2\norm{S}_{L_p}
   \le\beta^7\norm{S}_{L_p} .
\]
Finally $\log2>\tfrac12$.
\end{proof}

Taking bilinear suprema over deterministic unit vectors on a fixed
deleted submatrix, \eqref{mi32-real-doubling-bound} gives
\begin{equation}\label{mi32-weak-doubling}
 R_{X,I}(2p)\le J_\alpha R_{X,I}(p),\qquad p\ge\log2,
\end{equation}
for every deterministic $I$, since each bilinear form
$\sum_{i,j\notin I}X_{ij}s_it_j$ is a deterministic linear combination
of the original entries. The set $I$ is unchanged by this operation.
\section{The exact parity--amplitude representation}
\label{mi32-representation-section}

For now let $Y=(Y_{ij})_{i,j\le n}$ have independent symmetric entries
satisfying \eqref{mi32-regularity}. There is one original label
$e=(i,j)$ for each nonzero entry; denote this finite label set by
$\mathcal E$. Keeping zero entries as zero multipliers would give the
same formulas. Row vertices $r_i$ and column vertices $c_j$ belong to
disjoint copies of $[n]$; write $\mathcal V$ for their union.

Represent the joint law as $Y_e=\varepsilon_eZ_e$. The magnitudes are
independent because the original entries are independent. Signs can be
chosen independent of all magnitudes even at zero atoms: the product
law on $(\varepsilon,Z)$ has exactly the required image law for $Y$.
The sign at $(i,j)$ is distinct from that at $(j,i)$ when both labels
exist. Both directions of the single edge between $r_i$ and $c_j$
use the same variable $Y_{ij}$.

\subsection{The space and its polynomial domain}

Let $\mu_Z$ be the joint magnitude law and define
\begin{equation}\label{mi32-state-space}
 \mathcal H=L_2\bigl(\mu_Z;
                    \ell_2(\mathcal V\times2^{\mathcal E})\bigr),
 \qquad
 \norm{f}_{\mathcal H}^2
      =\E_Z\sum_{v\in\mathcal V}\sum_{S\subset\mathcal E}
                                                |f_{v,S}(Z)|^2.
\end{equation}
Conditional Fourier transform in the signs is the isometry
\begin{equation}\label{mi32-fourier-transform}
 (\mathcal U f)_v(\varepsilon,Z)
          =\sum_{S\subset\mathcal E}\varepsilon_S f_{v,S}(Z)
\end{equation}
onto $L_2(\varepsilon,Z;\ell_2(\mathcal V))$.
This transform preserves the $L_2$ norm, not general $L_p$ norms.
For example, with one sign and unit magnitude, amplitudes
$f_{\varnothing}=f_{\{e\}}=1$ have coefficient-vector norm $\sqrt2$,
whereas $\mathcal U f=1+\varepsilon_e$ has $L_4$ norm $2^{3/4}$.
Every non-$L_2$ norm used in the contraction below is therefore a norm
of the reconstructed random polynomial $\mathcal U f$, rather than
of its occupation-amplitude vector.
All formulas below may be restricted to polynomial coefficients, for
which every required moment is finite. No assertion about boundedness
of the full multiplication operator on $\mathcal H$ is needed.

Define $\mathcal P_d^+$ to be the image under $\mathcal U^{-1}$ of
vector polynomials in the original $Y_e$ of degree at most $d$ whose
ordinary coefficients are nonnegative in each current-vertex coordinate.
Thus a monomial $a_{v,\nu}Y^\nu$ contributes
$a_{v,\nu}Z^\nu$ to the component $(v,S)$ with
$S=\{e:\nu_e\text{ is odd}\}$. This parity compatibility is part
of the definition. It ensures that $\mathcal U f$ is a polynomial in
the original variables, so Lemma~\ref{mi32-cone-lemma} applies.

The parity grade of $(v,S)$ is $|S|$. Let $P_k$ keep exactly grade $k$.
For $f\in\mathcal P_d^+$, necessarily $P_kf=0$ for $k>d$.
Projection onto any collection of pairs $(v,S)$ simply discards
ordinary monomials, hence preserves $\mathcal P_d^+$. Distinct such
collections are orthogonal in \eqref{mi32-state-space}.
Magnitude coefficients on different parity states may have common
even powers, including powers of edges outside the occupied set $S$.
They are not required to be orthogonal or independent.

\subsection{Multiplication, creation, and annihilation}

The actual symmetric dilation is
\begin{equation}\label{mi32-dilation}
 H_Y=\begin{pmatrix}0&Y\\Y^\top&0\end{pmatrix}.
\end{equation}
For a scalar magnitude polynomial $h(Z)$ at state $(v,S)$,
conjugating multiplication by $H_Y$ through $\mathcal U$ gives
\begin{equation}\label{mi32-exact-operator}
 \mathcal T\bigl(h(Z)|v,S\rangle\bigr)
    =\sum_{e=vw}Z_eh(Z)|w,S\triangle\{e\}\rangle.
\end{equation}
Here $e=vw$ means the original bipartite entry with endpoints $v,w$.
Define $\mathcal C_+$ by keeping terms with $e\notin S$ and
$\mathcal C_-$ by keeping terms with $e\in S$. Their sum is exactly
$\mathcal T$, and they change parity grade by $+1$ and $-1$.

Both maps multiply the magnitude coefficient by $Z_e$. On an ordinary
monomial $Y^\nu$, either action increases $\nu_e$ by one. Therefore
\begin{equation}\label{mi32-cone-preservation}
 \mathcal C_+\mathcal P_d^+\subset\mathcal P_{d+1}^+,
 \qquad
 \mathcal C_-\mathcal P_d^+\subset\mathcal P_{d+1}^+,
 \qquad
 \mathcal T\mathcal P_d^+\subset\mathcal P_{d+1}^+.
\end{equation}
In particular annihilation does not divide by a magnitude or reset its
earlier powers. On $Y_e^3$, its action gives $Y_e^4$ while removing the
odd parity flag. This retains the full fourth moment of that same entry.

\section{Exact localization by the two original count vectors}
\label{mi32-counts-section}

For $S\subset\mathcal E$, define
\[
 r_i(S)=|S\cap(\{i\}\times[n])|,
 \qquad c_j(S)=|S\cap([n]\times\{j\})|.
\]
The vectors $r(S),c(S)$ count occupied original parity labels, not
total magnitude degrees. Write $\delta_i$ for a coordinate vector in
the relevant row or column count space.

\begin{lemma}[Two exact count decompositions]
\label{mi32-counts-lemma}
Consider maps from a current row to a current column, at input parity
grade $k$. Each of the following gives orthogonal partitions of both
the input states and the output states of the specified map.
\begin{enumerate}
\item For creation, label input $(r_i,S)$ by $d=r(S)+\delta_i$ and
output $(c_j,U)$ by $d=r(U)$. Put $I=\{i:d_i>0\}$. Then
$|I|\le k+1$ and the entire sector is an input/output parity-count
compression of multiplication by $Y[I,:]^\top$.
\item For annihilation, label input $(r_i,S)$ by $d=c(S)$ and output
$(c_j,U)$ by $d=c(U)+\delta_j$. Put $J=\{j:d_j>0\}$. Then
$|J|\le k$ and the entire sector is an input/output parity-count
compression of multiplication by $Y[:,J]^\top$. At $k=0$ the map is
zero.
\end{enumerate}
The reverse current-side direction transposes the statements.
Every magnitude variable remains available inside the coefficients.
\end{lemma}
\begin{proof}
Creation along $e=(i,j)\notin S$ produces
$U=S\cup\{(i,j)\}$, so $r(U)=r(S)+\delta_i=d$.
For fixed $d$, every input current row lies in $I$, and all parity
flags on input and output lie in $I\times[n]$.
The sum of the coordinates of $d$ is $k+1$, giving $|I|\le k+1$.
Multiplication by $Y[I,:]^\top$ therefore includes every actual
transition, with exactly its original coefficient.

More explicitly, let $P_d$ select the input label $r(S)+\delta_i=d$
and $Q_d$ select output grade $k+1$ with $r(U)=d$. The sector map is
$Q_d\mathcal U^{-1}M_{Y[I,:]^\top}\mathcal U P_d$, with the natural
current-coordinate inclusions understood. Output grade removes every
annihilation term. Conversely every selected creation term satisfies
the same original input and output conditions. This proves equality
of maps, rather than a comparison of their coefficients.

For annihilation, $e=(i,j)\in S$ produces
$U=S\setminus\{(i,j)\}$, so $c(U)+\delta_j=c(S)=d$.
Every parity flag lies in $[n]\times J$ and every output current
column lies in $J$, although the input current row may be arbitrary.
Now $\sum_jd_j=k$. Select input label $c(S)=d$ and output grade
$k-1$ with $c(U)+\delta_j=d$. The analogous compression of
$M_{Y[:,J]^\top}$ retains exactly the annihilation terms.
Creation terms have grade $k+1$ and cannot survive the output projection.

Each input label is a function of $(v,S)$ alone, and so is each output
label. Distinct labels thus give disjoint components of
\eqref{mi32-state-space}, hence orthogonal decompositions. The two
current-side directions have orthogonal domains and ranges. Different
input grades also have different output grades within each of
$\mathcal C_+$ and $\mathcal C_-$. No amplitude restriction was
used: even powers outside the selected row or column set remain in
$f_{v,S}(Z)$ throughout.
\end{proof}

The distinction between the two labels is necessary. The creator is
controlled through a small set of original rows in its output counts;
the annihilator is controlled through a small set of original columns
in its input counts. We will estimate the latter directly, not infer
it by applying an adjoint argument to a cone-restricted norm inequality.

\section{Thin matrices and the cone contraction}
\label{mi32-thin-section}

\begin{lemma}[A net only on the small axis]
\label{mi32-thin-lemma}
Let $q\ge1$ be an integer. For a fixed original row set $I$ and a fixed
original column set $J$,
\begin{align}
 \bigl(\E\norm{Y[I,:]}^{4q}\bigr)^{1/(4q)}
 &\le2\cdot5^{|I|/(4q)}
        \bigl\{\sigma_r(Y)+6\alpha^4R_Y(2q)\bigr\},
                     \label{mi32-thin-rows}\\
 \bigl(\E\norm{Y[:,J]}^{4q}\bigr)^{1/(4q)}
 &\le2\cdot5^{|J|/(4q)}
        \bigl\{\sigma_c(Y)+6\alpha^4R_Y(2q)\bigr\}.
                     \label{mi32-thin-columns}
\end{align}
In particular the numerical factor is at most $2\cdot5^{1/4}$ as soon as
the small axis has at most $q$ indices. Empty submatrices give zero and
require no net.
\end{lemma}
\begin{proof}
Fix a deterministic unit vector $s\in\R^I$ and consider the random
vector $Y[I,:]^\top s$, whose $j$-th coordinate is
\[
 V_j=\sum_{i\in I}s_iY_{ij}.
\]
Distinct coordinates use disjoint sets of original entries, so the
family $(V_j)_j$ is independent; each $V_j$ is a deterministic linear
combination of independent symmetric variables, hence symmetric in law.
We verify the three hypotheses of
Theorem~\ref{mi32-hilbert-comparison} for this family.

First, independence and zero means give
\begin{equation}\label{mi32-fixed-vector-variance}
 \sigma^2=\sum_j\E V_j^2
   =\sum_{i\in I}s_i^2\sum_j\E Y_{ij}^2
   \le\sigma_r(Y)^2 .
\end{equation}
Second, Corollary~\ref{mi32-scalar-doubling} applies to each $V_j$ and
supplies \eqref{mi32-coordinate-regularity} with
$\beta=H_\alpha$. Third, for a deterministic unit vector $t$,
\[
 \sum_jt_jV_j=\sum_{i\in I,\,j}s_it_jY_{ij},
\]
which is one of the bilinear forms in the definition
\eqref{mi32-weak} of $R_Y$ evaluated at the unit vectors
$s$ (extended by zero) and $t$; hence $W_{2q}\le R_Y(2q)$.
Theorem~\ref{mi32-hilbert-comparison} therefore gives
\begin{equation}\label{mi32-fixed-vector-strong}
 \bigl\|\,\norm{Y[I,:]^\top s}_2\,\bigr\|_{L_{4q}}
 \le\sigma_r(Y)+2H_\alpha^2R_Y(2q)
  =\sigma_r(Y)+6\alpha^4R_Y(2q),
\end{equation}
because $2H_\alpha^2=2\cdot3\alpha^4=6\alpha^4$.

Choose a deterministic $1/2$-net $\mathcal N$ of the unit sphere in
$\R^I$, with $|\mathcal N|\le5^{|I|}$; this cardinality follows by
taking a maximal $1/2$-separated set and comparing its disjoint balls
of radius $1/4$ with the ball of radius $5/4$. For every realization,
$\norm{Y[I,:]}\le2\max_{s\in\mathcal N}\norm{Y[I,:]^\top s}_2$.
Raising to the power $4q$ and bounding a maximum by a sum,
\[
 \bigl(\E\norm{Y[I,:]}^{4q}\bigr)^{1/(4q)}
 \le2\Bigl(\sum_{s\in\mathcal N}
        \E\norm{Y[I,:]^\top s}_2^{4q}\Bigr)^{1/(4q)}
 \le2\cdot5^{|I|/(4q)}
        \bigl\{\sigma_r(Y)+6\alpha^4R_Y(2q)\bigr\},
\]
which is \eqref{mi32-thin-rows}. Transposing the matrix exchanges the
roles of rows and columns and proves \eqref{mi32-thin-columns} with the
maximum column variance. If $|I|\le q$ then $5^{|I|/(4q)}\le5^{1/4}$.
\end{proof}

The weak moment entering \eqref{mi32-thin-rows} is taken at order
$2q$, half the order of the strong moment on the left. This is a feature
of Theorem~\ref{mi32-hilbert-comparison} and is what makes the whole
chain close at the logarithmic order: the local theorem below needs its
weak input at $2q$, not at $4q$.

\begin{proposition}[Contraction on the preserved cone]
\label{mi32-contraction-proposition}
For an integer $q\ge1$, put
\begin{equation}\label{mi32-B-constant}
 B_\alpha=2\cdot5^{1/4}H_\alpha=2\sqrt3\cdot5^{1/4}\alpha^2 .
\end{equation}
For every $f\in\mathcal P_{q-1}^+$,
\begin{equation}\label{mi32-contraction}
 \begin{aligned}
 \norm{\mathcal C_\pm f}_{\mathcal H}
   &\le B_\alpha\bigl\{\sigma(Y)+6\alpha^4R_Y(2q)\bigr\}
            \norm{f}_{\mathcal H},\\
 \norm{\mathcal T f}_{\mathcal H}
   &\le2B_\alpha\bigl\{\sigma(Y)+6\alpha^4R_Y(2q)\bigr\}
            \norm{f}_{\mathcal H}.
 \end{aligned}
\end{equation}
\end{proposition}
\begin{proof}
Set $p=4q$. First consider row-to-column creation at a fixed parity
grade $k\le q-1$ and decompose its input as $\sum_df_d$ using
Lemma~\ref{mi32-counts-lemma}. Every $f_d$ lies in
$\mathcal P_{q-1}^+$, since its projection only discards ordinary
monomials. Its sector has $|I|\le k+1\le q$.
Apply $\mathcal U$ first. It identifies the multiplier in that sector
with $Y[I,:]^\top$, and the transformed output projection is an
$L_2$ contraction. H\"older is then applied on the original joint
probability space to $Y[I,:]^\top\mathcal U f_d$. Together,
Lemma~\ref{mi32-cone-lemma} and
Lemma~\ref{mi32-thin-lemma} imply
\begin{align}
 \norm{\mathcal C_{+,d}f_d}_{\mathcal H}
 &\le\norm{Y[I,:]^\top\mathcal U f_d}_{L_2(\ell_2)}\nonumber\\
 &\le\bigl(\E\norm{Y[I,:]}^p\bigr)^{1/p}
            \norm{\mathcal U f_d}_{L_{2p/(p-2)}(\ell_2)}\nonumber\\
 &\le2\cdot5^{1/4}\bigl\{\sigma_r(Y)+6\alpha^4R_Y(2q)\bigr\}
      H_\alpha^{4(q-1)/p}\norm{f_d}_{\mathcal H}\nonumber\\
 &\le B_\alpha\bigl\{\sigma_r(Y)+6\alpha^4R_Y(2q)\bigr\}
            \norm{f_d}_{\mathcal H}.
 \label{mi32-sector-bound}
\end{align}
In the last step $4(q-1)/p=(q-1)/q\le1$ and $H_\alpha\ge1$.
The expectation in H\"older is over the full joint sign and magnitude
law. It requires no independence between the matrix multiplier and
$\mathcal U f_d$, so magnitude powers outside $I$ stay in the state.
The $L_{2p/(p-2)}$ factor is explicitly the norm of $\mathcal U f_d$.
Only its final $L_2$ norm is replaced by $\norm{f_d}_{\mathcal H}$
using the isometry; no $L_p$ isometry of occupation amplitudes is used.

For row-to-column annihilation use its column-count sectors instead.
They have $|J|\le k\le q-1$, and \eqref{mi32-thin-columns} gives the
same bound with $\sigma_c(Y)$. The reverse current-side direction
transposes these two statements. Sum the squared estimates over the
orthogonal input and output sectors, over their grades, and over
current-side directions. The sums of squared input norms equal
$\norm{f}_{\mathcal H}^2$; there is no factor for the number of
sectors. This proves the separate bounds for $\mathcal C_+$ and
$\mathcal C_-$, after replacing $\sigma_r$ and $\sigma_c$ by their
maximum $\sigma$. Their triangle inequality gives the stated bound for
$\mathcal T$.
\end{proof}

There is no random maximum over row or column sets in this argument.
Each set is the support of a deterministic count label, and the same
unconditional inequality bounds every corresponding squared sector
norm before summation. Nor is \eqref{mi32-contraction} a bound on the
entire polynomial linear span: it is asserted only on the specified
cone at the specified total degree.

\section{Vacuum walks and the local moment theorem}
\label{mi32-vacuum-section}

\begin{theorem}[Symmetric regular-law local moment bound]
\label{mi32-local-theorem}
Let $Y$ be an $n\times n$ matrix with independent symmetric entries
satisfying \eqref{mi32-regularity}. With $B_\alpha$ from
\eqref{mi32-B-constant}, every integer $q\ge1$ satisfies
\begin{equation}\label{mi32-local-moment}
 \bigl(\E\norm{Y}^{2q}\bigr)^{1/(2q)}
 \le(2n)^{1/(2q)}\,2B_\alpha
       \bigl\{\sigma(Y)+6\alpha^4R_Y(2q)\bigr\}.
\end{equation}
The same bound holds with $\sigma(Y)$ replaced by any $B$ that
dominates both the row and the column variance scale, in particular by
$M(Y)$.
\end{theorem}
\begin{proof}
For a current vertex $v\in\mathcal V$, let $\Omega_v$ be the state
with coefficient one at $(v,\varnothing)$ and zero elsewhere. It has
norm one in $\mathcal H$. Under $\mathcal U$ it is the constant vector
$e_v$. For every $j\ge0$,
\[
 \mathcal U\mathcal T^j\Omega_v=H_Y^j e_v.
\]
Each coordinate on the right is the sum, over original ordered walks
of length $j$ between the indicated vertices, of the product of their
entry variables. Every coefficient is a nonnegative integer, even
when the realized entries are negative. Repeated edges retain their
powers in the same monomial. Thus
$\mathcal T^j\Omega_v\in\mathcal P_j^+$.

Apply Proposition~\ref{mi32-contraction-proposition} successively to
the prefixes $j=0,\ldots,q-1$, keeping the common horizon $q$:
\begin{equation}\label{mi32-vacuum-prefix}
 \norm{\mathcal T^q\Omega_v}_{\mathcal H}
 \le\bigl[2B_\alpha\{\sigma(Y)+6\alpha^4R_Y(2q)\}\bigr]^q .
\end{equation}
The dilation $H_Y$ is real symmetric, so $\norm{H_Y}^q=\norm{H_Y^q}$
and, comparing the operator norm of $H_Y^q$ with its Frobenius norm,
\begin{equation}\label{mi32-vacuum-trace}
 \norm{H_Y}^{2q}\le\sum_{v\in\mathcal V}\norm{H_Y^qe_v}_2^2 .
\end{equation}
Taking expectations, using
$\norm{\mathcal T^q\Omega_v}_{\mathcal H}^2
 =\E\norm{H_Y^qe_v}_2^2$ and $|\mathcal V|=2n$,
\[
 \E\norm{H_Y}^{2q}
 \le2n\bigl[2B_\alpha\{\sigma(Y)+6\alpha^4R_Y(2q)\}\bigr]^{2q}.
\]
Finally $\norm{Y}\le\norm{H_Y}$, because $H_Y$ maps the embedded copy
of a column vector to the embedded copy of its image under $Y$
isometrically on each side. Taking the $2q$-th root gives
\eqref{mi32-local-moment}. Only the row and column variance budgets
enter Proposition~\ref{mi32-contraction-proposition}, so any common
bound $B$ for both may replace $\sigma(Y)$.
\end{proof}

The dimension factor counts both sides of the bipartite dilation. One
may instead compare the two diagonal blocks of the even power and
replace $2n$ by $n$; we retain $2n$ because it is the factor the formal
development certifies, and at $q$ comparable to $\log n$ the difference
is at most $\sqrt2$. Every equality above uses the original joint law,
and all expectations are finite by the entry moment assumptions. The
proof requires neither polynomial completeness in an infinite $L_2$
space nor a bounded extension of the full multiplication operator.

\begin{corollary}[The local mean at the exact logarithmic order]
\label{mi32-local-mean}
Let $Y$ be as in Theorem~\ref{mi32-local-theorem} with $n\ge1$, and put
\begin{equation}\label{mi32-L-constant}
 L_\alpha=4e\cdot5^{1/4}\sqrt3\,\alpha^2=2e\,B_\alpha .
\end{equation}
Then
\begin{equation}\label{mi32-local-mean-bound}
 \E\norm{Y}\le L_\alpha
    \bigl\{M(Y)+6\alpha^4J_\alpha^3R_Y(\log(n+1))\bigr\}.
\end{equation}
\end{corollary}
\begin{proof}
Take $q=\lceil2\log(n+1)\rceil\ge1$. Then
$\log(2n)\le2\log(n+1)\le q$, so $(2n)^{1/(2q)}\le e^{1/2}\le e$, and
$2q\le4\log(n+1)+2\le8\log(n+1)$ because
$\log(n+1)\ge\log2>\tfrac12$. Three applications of
\eqref{mi32-weak-doubling}, with $I=\varnothing$, therefore give
$R_Y(2q)\le J_\alpha^3R_Y(\log(n+1))$. Since
$\E\norm{Y}\le(\E\norm{Y}^{2q})^{1/(2q)}$, the claim follows from
\eqref{mi32-local-moment} with $B=M(Y)\ge\sigma(Y)$.
\end{proof}

The right-hand side of \eqref{mi32-local-mean-bound} uses the
\emph{undeleted} weak moment $R_Y=R_{Y,\varnothing}$. Replacing it by
the smaller deletion parameter $D(Y)$ is exactly the remaining matrix
problem, and it is the subject of the next section.
\section{From the local estimate to shared-index deletion}
\label{mi32-reduction-section}

Corollary~\ref{mi32-local-mean} controls $\E\norm{Y}$ by the undeleted
weak moment $R_Y(\log(n+1))$, which can be much larger than $D(Y)$.
This section closes that gap. The target is the symmetric-law statement

\begin{theorem}[Symmetric-law deletion bound]
\label{mi32-symmetric-deletion}
There is a constant $C^{\rm sym}_\alpha$, depending only on $\alpha$,
such that every $n\ge1$ and every $n\times n$ matrix $Y$ with
independent entries that are symmetric in law and satisfy
\eqref{mi32-regularity} obey
\begin{equation}\label{mi32-symmetric-deletion-bound}
 \E\norm{Y}\le C^{\rm sym}_\alpha\{M(Y)+D(Y)\}.
\end{equation}
One may take
\begin{equation}\label{mi32-Csym}
 C^{\rm sym}_\alpha
 =\sqrt{31}+2\sqrt{1+\sqrt2\,\alpha^2}
   +24e\cdot5^{1/4}\sqrt3\,\alpha^2
   +144e\cdot5^{1/4}\sqrt3\,\alpha^6J_\alpha^{10}.
\end{equation}
\end{theorem}

Section~\ref{mi32-symmetrization} then removes the symmetry assumption.
The overall design of the reduction --- a nested deterministic selection
of indices, a decomposition into near bands and a far remainder, and a
maximum over blocks with growing moment orders --- follows
Lata\l a--\'Swi\k atkowski\ \cite[Remark~4.5]{LatalaSwiatkowski}. Two
steps are different here, and both differences are what make the
argument self-contained.

\begin{enumerate}
\item The published remainder step invokes a variance-envelope estimate
that is a consequence of
\cite[Theorem~4.4]{LatalaVanHandelYoussef}. Adapting that theorem to
the present hypotheses would require a one-sided version of a Gaussian
comparison whose stated form assumes two-sided polynomial moment
growth, which $\alpha$-regular entries need not have. We instead raise
the per-coordinate screening threshold from $b^2$ to $b^3$ and estimate
the remainder by an elementary centered Gram second moment
(Section~\ref{mi32-far-remainder}). No general fourth-moment matrix
inequality is used.
\item The published argument works with a matrix-symmetric dilation,
whose reflected entries repeat variables, and afterwards pays for a
paired deletion budget. We instead run the whole decomposition on the
original rectangular matrix with one shared index family on both axes
--- which is precisely what the functional $D$ in \eqref{mi32-D}
quantifies over. Every entry family used below is therefore genuinely
independent, and no dilation budget arises.
\end{enumerate}

Throughout this section $Y$ is as in
Theorem~\ref{mi32-symmetric-deletion}, and we abbreviate
\[
 M=M(Y),\qquad D=D(Y),\qquad v_{ij}=\E Y_{ij}^2 .
\]
By \eqref{mi32-variances}, $\sum_jv_{ij}\le M^2$ for every $i$ and
$\sum_iv_{ij}\le M^2$ for every $j$. Regularity at order two gives the
fourth-moment budget
\begin{equation}\label{mi32-fourth-budget}
 \E Y_{ij}^4\le\alpha^4v_{ij}^2 ,
\end{equation}
since $\norm{Y_{ij}}_{L_4}\le\alpha\norm{Y_{ij}}_{L_2}$.

\subsection{The deletion schedule}
\label{mi32-schedule}

Define, for $k\ge0$,
\begin{equation}\label{mi32-schedule-def}
 b_k=2^{5^k}-1,\qquad
 \pi_k=\log(b_k+1)=5^k\log2,\qquad
 t_k=b_k^3 .
\end{equation}
So $b_0=1$, $\pi_0=\log2$, $b_1=31$, and $b_{k+1}+1=(b_k+1)^5$.
Three elementary facts are used.

\begin{lemma}[Schedule arithmetic]
\label{mi32-schedule-lemma}
For every $k\ge0$ and every integer $0\le a\le b_k$,
\begin{align}
 2\bigl(a+a\,t_k\bigr)+b_k&\le b_{k+1},
                     \label{mi32-schedule-card}\\
 2^k&\le b_k,\qquad n\le b_n\quad(n\ge0),
                     \label{mi32-schedule-growth}\\
 5^{a+1}&\le125\cdot5^{(a-2)^+}\quad(a\ge1).
                     \label{mi32-schedule-ratio}
\end{align}
\end{lemma}
\begin{proof}
Write $b=b_k\ge1$. Expanding $b_{k+1}=(b+1)^5-1$ gives
$b_{k+1}=b^5+5b^4+10b^3+10b^2+5b$, whereas
$2(a+ab^3)+b\le2(b+b^4)+b=2b^4+3b$; since $b\ge1$ this is at most
$5b^4+5b\le b_{k+1}$, proving \eqref{mi32-schedule-card}. For
\eqref{mi32-schedule-growth}, $b_k+1=2^{5^k}\ge2^{k+1}$ because
$5^k\ge k+1$, and $n+1\le2^n\le2^{5^n}=b_n+1$. Finally
\eqref{mi32-schedule-ratio} is an equality for $a\ge2$ and reads
$25\le125$ at $a=1$.
\end{proof}

\subsection{Two-sided screening and the nested family}
\label{mi32-nested}

Fix a deterministic index set $S$ and an integer $t\ge0$. For a column
$j\in S$, at most $t$ rows $i$ can satisfy $(t+1)v_{ij}>M^2$: otherwise
$t+1$ of them would already contribute more than $M^2$ to
$\sum_iv_{ij}\le M^2$. Let $\mathrm{col}_t(S)$ be the union of $S$ with
all such rows, over all $j\in S$, and let $\mathrm{row}_t(S)$ be the
transposed construction. Then
\begin{equation}\label{mi32-selection-card}
 |\mathrm{col}_t(S)|\le|S|(1+t),\qquad
 |\mathrm{row}_t(S)|\le|S|(1+t),
\end{equation}
and, by construction,
\begin{equation}\label{mi32-selection-screen}
 j\in S,\ i\notin\mathrm{col}_t(S)\ \Longrightarrow\
 (t+1)v_{ij}\le M^2,
\end{equation}
with the analogous statement for $\mathrm{row}_t$ and $v_{ji}$.

For an integer budget $k\ge1$ let $I_k$ be a deterministic minimizer of
$I\mapsto R_{Y,I}(\log(k+1))$ over $|I|\le k$; the minimum is attained
because there are finitely many sets, so
\begin{equation}\label{mi32-minimizer}
 R_{Y,I_k}(\log(k+1))=\min_{|I|\le k}R_{Y,I}(\log(k+1))\le D
 \qquad(1\le k\le n),
\end{equation}
directly from \eqref{mi32-D}. For $k>n$ we may take $I_k=[n]$, and then
$R_{Y,I_k}(p)=0$ for every $p>0$, so \eqref{mi32-minimizer} holds for
every $k\ge1$. Now define the nested family
\begin{equation}\label{mi32-nested-def}
 S_0=I_{b_0}=I_1,\qquad
 S_{k+1}=\mathrm{col}_{t_k}(S_k)\cup\mathrm{row}_{t_k}(S_k)
                \cup I_{b_k}.
\end{equation}

\begin{lemma}[The nested family]
\label{mi32-nested-lemma}
The sets \eqref{mi32-nested-def} are increasing, and for every $k\ge0$
\begin{equation}\label{mi32-nested-card}
 |S_k|\le b_k,\qquad S_{n+1}=[n].
\end{equation}
Moreover, for $r+1\le s$,
\begin{equation}\label{mi32-nested-screen}
 j\in S_r,\ i\notin S_s\ \Longrightarrow\
 v_{ij}\le\frac{M^2}{b_r^3+1},
 \qquad
 i\in S_r,\ j\notin S_s\ \Longrightarrow\
 v_{ij}\le\frac{M^2}{b_r^3+1},
\end{equation}
and, for every $a\ge1$,
\begin{equation}\label{mi32-nested-weak}
 R_{Y,S_{a-1}}\bigl(\pi_{(a-2)^+}\bigr)\le D .
\end{equation}
\end{lemma}
\begin{proof}
Monotonicity is clear from $S\subset\mathrm{col}_t(S)$. Since
$|S_0|\le b_0$, \eqref{mi32-selection-card} and
\eqref{mi32-schedule-card} give
$|S_{k+1}|\le2|S_k|(1+t_k)+b_k\le b_{k+1}$ inductively. By
\eqref{mi32-schedule-growth}, $b_n\ge n$, so $I_{b_n}=[n]$ may be
taken and $S_{n+1}=[n]$.

For \eqref{mi32-nested-screen}, note $i\notin S_s$ and $r+1\le s$ imply
$i\notin S_{r+1}$, hence $i\notin\mathrm{col}_{t_r}(S_r)$; now
\eqref{mi32-selection-screen} with $t=t_r=b_r^3$ gives the first
implication, and the transposed construction gives the second.

For \eqref{mi32-nested-weak}, enlarging a deleted set can only decrease
the weak moment, because masking both deterministic test vectors on the
larger set is a Euclidean contraction and reproduces exactly the
smaller bilinear form. At $a=1$ we have $S_{a-1}=S_0=I_1$ and
$\pi_{(a-2)^+}=\pi_0=\log2=\log(1+1)$, so \eqref{mi32-minimizer} at
$k=1$ applies. At $a\ge2$ we have $I_{b_{a-2}}\subset S_{a-1}$ by
\eqref{mi32-nested-def}, and $\pi_{a-2}=\log(b_{a-2}+1)$, so
\[
 R_{Y,S_{a-1}}(\pi_{a-2})\le R_{Y,I_{b_{a-2}}}(\pi_{a-2})\le D
\]
by \eqref{mi32-minimizer} at $k=b_{a-2}$.
\end{proof}

Only $b_r$, and never $b_{r+1}$, appears in \eqref{mi32-nested-screen}
and \eqref{mi32-nested-weak}: the minimizer inserted at step $k$ carries
the budget $b_k$, one stage behind the set that contains it. This
one-stage shift is what keeps the recorded logarithmic orders
comparable to the scheduled ones in Section~\ref{mi32-near-bands}.

\subsection{Shells and the exact five-mask partition}
\label{mi32-shells}

By \eqref{mi32-nested-card} every index lies in some $S_k$, so
\begin{equation}\label{mi32-level-def}
 \ell_i=\min\{r\ge0:\ i\in S_r\}\in\{0,1,\ldots,n+1\}
\end{equation}
is well defined, and $\{i:\ell_i\le r\}=S_r$ by monotonicity. Hence
\begin{equation}\label{mi32-level-card}
 \bigl|\{i:\ \ell_i\le r\}\bigr|\le b_r .
\end{equation}
Split the entries of $Y$ by the five deterministic masks
\begin{equation}\label{mi32-masks}
 \begin{aligned}
 \text{\rm nearLow}&:\ \ell_i\le1\ \text{and}\ \ell_j\le1, &
 \text{\rm farLow}&:\ \ell_j+1<\ell_i,\\
 \text{\rm nearHigh}&:\ \lfloor\ell_i/2\rfloor
     =\lfloor\ell_j/2\rfloor\ge1, &
 \text{\rm farUp}&:\ \ell_i+1<\ell_j,\\
 \text{\rm nearCross}&:\ \bigl\lfloor(\ell_i+1)/2\bigr\rfloor
     =\bigl\lfloor(\ell_j+1)/2\bigr\rfloor
        \ \text{and}\ \ell_i\ne\ell_j. &&
 \end{aligned}
\end{equation}

\begin{lemma}[Exact partition]
\label{mi32-partition-lemma}
For all integers $a,b\ge0$, exactly one of the five conditions in
\eqref{mi32-masks} holds at $(\ell_i,\ell_j)=(a,b)$. Consequently
$Y$ is the sum of the five corresponding masked matrices, and
\begin{equation}\label{mi32-five-masks}
 \norm{Y}\le\norm{Y^{\rm nL}}+\norm{Y^{\rm nH}}
          +\norm{Y^{\rm nC}}+\norm{Y^{\rm fL}}+\norm{Y^{\rm fU}}
\end{equation}
pointwise, where the superscripts name the masks in the order listed.
\end{lemma}
\begin{proof}
The two far conditions are disjoint from each other and from
$|a-b|\le1$. For $|a-b|\le1$: if $\lfloor a/2\rfloor=\lfloor
b/2\rfloor$ then either both $a,b\le1$, which is nearLow, or
$\lfloor a/2\rfloor\ge1$, which is nearHigh, and these two are
disjoint; otherwise $|a-b|=1$ with $a,b$ in different pairs
$\{2m,2m+1\}$, which forces $\{a,b\}=\{2m+1,2m+2\}$ and hence
$\lfloor(a+1)/2\rfloor=\lfloor(b+1)/2\rfloor=m+1$ with $a\ne b$, which
is nearCross and excludes $\lfloor a/2\rfloor=\lfloor b/2\rfloor$.
Conversely nearCross forces $|a-b|=1$. The triangle inequality gives
\eqref{mi32-five-masks}.
\end{proof}

\subsection{The far remainder has a summable Gram energy}
\label{mi32-far-remainder}

\begin{proposition}[Far orientations]
\label{mi32-far-proposition}
With $Y^{\rm fL}$ and $Y^{\rm fU}$ as in
Lemma~\ref{mi32-partition-lemma},
\begin{equation}\label{mi32-far-bound}
 \E\norm{Y^{\rm fL}}\le M\sqrt{1+\sqrt2\,\alpha^2},
 \qquad
 \E\norm{Y^{\rm fU}}\le M\sqrt{1+\sqrt2\,\alpha^2}.
\end{equation}
\end{proposition}
\begin{proof}
Write $A=Y^{\rm fL}$, so $A_{ij}=Y_{ij}\mathbf1\{\ell_j+1<\ell_i\}$,
and put $w_{ij}=\E A_{ij}^2\le v_{ij}$. The masks are deterministic, so
the entries of $A$ are still independent, centered and symmetric, and
\eqref{mi32-fourth-budget} still holds for them. Also
$\sum_iw_{ij}\le M^2$ for every $j$.

\emph{Support cap.} Suppose $w_{ik}\ne0$ and $\ell_j\le\ell_k$. Then
$\ell_k+1<\ell_i$, so $i\notin S_{\ell_k+1}$, while $j\in
S_{\ell_k}$. Applying the first implication of
\eqref{mi32-nested-screen} with $r=\ell_k$ and $s=\ell_k+1$,
\begin{equation}\label{mi32-far-cap}
 w_{ij}\le v_{ij}\le\frac{M^2}{b_{\ell_k}^3+1}.
\end{equation}

\emph{Pair energy.} Fix columns $j,k$ and let $r=\max(\ell_j,\ell_k)$.
By symmetry of the claim assume $\ell_j\le\ell_k$, so $r=\ell_k$. Only
indices $i$ with $w_{ik}\ne0$ contribute, and for those
\eqref{mi32-far-cap} applies, whence
\begin{equation}\label{mi32-far-pair}
 \sum_iw_{ij}w_{ik}
 \le\frac{M^2}{b_r^3+1}\sum_iw_{ik}
 \le\frac{M^4}{b_r^3+1}.
\end{equation}

\emph{Summation.} The number of ordered pairs $(j,k)$ with
$\max(\ell_j,\ell_k)=r$ is at most $|\{i:\ell_i\le r\}|^2\le b_r^2$ by
\eqref{mi32-level-card}. Since $b_r\ge1$ we have
$b_r^2/(b_r^3+1)\le1/b_r$, and $b_r\ge2^r$ by
\eqref{mi32-schedule-growth}. Therefore
\begin{equation}\label{mi32-far-total}
 \sum_j\sum_k\sum_iw_{ij}w_{ik}
 \le M^4\sum_{r\ge0}\frac{b_r^2}{b_r^3+1}
 \le M^4\sum_{r\ge0}2^{-r}
 =2M^4 .
\end{equation}

\emph{Centered Gram.} Let $d_j=\sum_iw_{ij}$, let
$\Delta=\diag(d_j)$, and put $G=A^\top A-\Delta$. Distinct original
entries are independent, and masking preserves this. For $j\ne k$,
\[
 G_{jk}=\sum_iA_{ij}A_{ik}
\]
is a sum over $i$ of independent centered products, so $\E G_{jk}=0$
and, the cross terms vanishing,
\begin{equation}\label{mi32-gram-off}
 \E G_{jk}^2=\sum_i\E\bigl(A_{ij}^2A_{ik}^2\bigr)
            =\sum_iw_{ij}w_{ik}.
\end{equation}
For $j=k$, $G_{jj}=\sum_i(A_{ij}^2-w_{ij})$ is a sum of independent
centered variables, so by \eqref{mi32-fourth-budget}
\begin{equation}\label{mi32-gram-diag}
 \E G_{jj}^2=\sum_i\bigl(\E A_{ij}^4-w_{ij}^2\bigr)
    \le\alpha^4\sum_iw_{ij}^2 .
\end{equation}
No independence among the entries of $G$ is asserted or needed. Since
$\alpha\ge1$, the right-hand sides of \eqref{mi32-gram-off} and
\eqref{mi32-gram-diag} are at most $\alpha^4$ times the corresponding
term of the triple sum in \eqref{mi32-far-total}, so
\begin{equation}\label{mi32-gram-frobenius}
 \E\norm{G}_F^2=\sum_j\sum_k\E G_{jk}^2\le2\alpha^4M^4 .
\end{equation}

\emph{From the Gram energy to the norm.} Deterministically
$\norm{A}^2=\norm{A^\top A}\le\norm{\Delta}+\norm{G}
 \le M^2+\norm{G}_F$, since
$\norm{\Delta}=\max_jd_j\le M^2$ and the operator norm is at most the
Frobenius norm. Two applications of Cauchy--Schwarz and
\eqref{mi32-gram-frobenius} give
\[
 \E\norm{A}\le\bigl(\E\norm{A}^2\bigr)^{1/2}
  \le\bigl(M^2+\E\norm{G}_F\bigr)^{1/2}
  \le\Bigl(M^2+\bigl(\E\norm{G}_F^2\bigr)^{1/2}\Bigr)^{1/2}
  \le M\sqrt{1+\sqrt2\,\alpha^2},
\]
which is the first bound in \eqref{mi32-far-bound}. All the
integrability used is that of the finitely many fourth entry moments.
There is no division by $M$, so $M=0$ is included.

For $Y^{\rm fU}$, note that its transpose is the lower-far mask of
$Y^\top$ for the same level function, and the operator norm is
invariant under transposition. Running the argument for $Y^\top$ uses
the second implication of \eqref{mi32-nested-screen} and the row
variance budget, and gives the same bound.
\end{proof}

\begin{remark}[Why the cube]
\label{mi32-cube-remark}
With the threshold $t_r=b_r^2$ the individual fourth moments are still
summable, but \eqref{mi32-far-total} would contribute
$b_r^2/(b_r^2+1)\approx1$ for every shell, and the number of shells is
not bounded in terms of $\alpha$. Recovering the estimate from that
route requires a general fourth-moment matrix inequality --- exactly the
external input we are avoiding. The cube absorbs the off-diagonal Gram
pairs directly, at no cost anywhere else: the enlarged selection is
charged only in \eqref{mi32-schedule-card}, which the schedule
\eqref{mi32-schedule-def} already accommodates.
\end{remark}

\subsection{The near bands}
\label{mi32-near-bands}

For $a\ge0$ put $V_a=\{i:\ \ell_i\in\{a,a+1\}\}$. Then
\begin{equation}\label{mi32-block-support}
 V_a\subset S_{a+1},\qquad |V_a|\le b_{a+1},
 \qquad\text{and}\qquad
 V_a\cap S_{a-1}=\varnothing\quad(a\ge1),
\end{equation}
the last because every $i\in V_a$ has $\ell_i\ge a$. Each of the two
remaining near masks is block diagonal:
$Y^{\rm nH}$ vanishes unless $\lfloor\ell_i/2\rfloor
=\lfloor\ell_j/2\rfloor$, and $Y^{\rm nC}$ vanishes unless
$\lfloor(\ell_i+1)/2\rfloor=\lfloor(\ell_j+1)/2\rfloor$. The block of
$Y^{\rm nH}$ at index $m$ is empty for $m=0$ and equals the full
submatrix $Y[V_{2m},V_{2m}]$ for $m\ge1$; the block of $Y^{\rm nC}$ at
index $m$ is empty for $m=0$ and, for $m\ge1$, is the part of
$Y[V_{2m-1},V_{2m-1}]$ supported off the two diagonal shell blocks.

\begin{lemma}[Block-diagonal norms and two-colour masks]
\label{mi32-block-norm-lemma}
Let $A$ be a matrix that vanishes whenever $c(i)\ne c(j)$ for a
deterministic map $c$. Then $\norm{A}=\max_m\norm{A[c^{-1}(m),
c^{-1}(m)]}$; in particular there is no factor for the number of
blocks. If moreover $\sigma$ takes two values on $c^{-1}(m)$ and
$A^{\rm off}$ agrees with $A$ where $\sigma(i)\ne\sigma(j)$ and
vanishes elsewhere, then $\norm{A^{\rm off}[c^{-1}(m),c^{-1}(m)]}
\le\norm{A[c^{-1}(m),c^{-1}(m)]}$.
\end{lemma}
\begin{proof}
The Euclidean space splits orthogonally along the fibres of $c$, and
$A$ preserves that splitting, which gives the first claim. For the
second, let $\Delta=\diag(\pm1)$ be $+1$ on one value of $\sigma$ and
$-1$ on the other. Then $A^{\rm off}=\tfrac12(A-\Delta A\Delta)$
entrywise on the block, and $\norm{\Delta}=1$, so
$\norm{A^{\rm off}}\le\tfrac12(\norm{A}+\norm{A})=\norm{A}$.
\end{proof}

\begin{proposition}[Blocks at the scheduled order]
\label{mi32-block-proposition}
For $a\ge1$ put $q_a=\lceil2\pi_{a+1}\rceil$ and
\begin{equation}\label{mi32-block-bound-constant}
 L=L_\alpha\bigl\{M+6\alpha^4J_\alpha^{10}D\bigr\},
\end{equation}
with $L_\alpha$ from \eqref{mi32-L-constant}. Then
\begin{equation}\label{mi32-block-bound}
 \bigl(\E\norm{Y[V_a,V_a]}^{2q_a}\bigr)^{1/(2q_a)}\le L,
 \qquad\text{and}\qquad
 a+2\le2q_a .
\end{equation}
\end{proposition}
\begin{proof}
The submatrix $Y[V_a,V_a]$ has independent symmetric entries satisfying
\eqref{mi32-regularity}, and its row and column variance sums are
sub-sums of those of $Y$, hence at most $M^2$.

\emph{Weak input.} Let $s,t$ be deterministic unit vectors supported on
$V_a$. Since $V_a\cap S_{a-1}=\varnothing$ by
\eqref{mi32-block-support}, the bilinear form $\sum_{i,j\in
V_a}Y_{ij}s_it_j$ is one of the forms occurring in
$R_{Y,S_{a-1}}$, so
$R_{Y[V_a,V_a]}(p)\le R_{Y,S_{a-1}}(p)$ for every $p>0$. Next,
\[
 2q_a\le4\pi_{a+1}+2
    \le500\,\pi_{(a-2)^+}+4\log2
    \le504\,\pi_{(a-2)^+}
    \le2^{10}\pi_{(a-2)^+},
\]
where the second step used \eqref{mi32-schedule-ratio} in the form
$4\cdot5^{a+1}\le500\cdot5^{(a-2)^+}$ and $2\le4\log2$, and the third
used $\pi_{(a-2)^+}\ge\log2$. Ten applications of
\eqref{mi32-weak-doubling} and then \eqref{mi32-nested-weak} give
\begin{equation}\label{mi32-block-weak}
 R_{Y[V_a,V_a]}(2q_a)
 \le J_\alpha^{10}R_{Y,S_{a-1}}\bigl(\pi_{(a-2)^+}\bigr)
 \le J_\alpha^{10}D .
\end{equation}

\emph{Dimension factor.} By \eqref{mi32-block-support},
$|V_a|\le b_{a+1}$, and $2b\le(b+1)^2$ for $b\ge1$ gives
\[
 \log\bigl(2|V_a|\bigr)\le\log\bigl(2b_{a+1}\bigr)
 \le2\log(b_{a+1}+1)=2\pi_{a+1}\le q_a\le2q_a ,
\]
so $(2|V_a|)^{1/(2q_a)}\le e$. Applying
Theorem~\ref{mi32-local-theorem} to $Y[V_a,V_a]$ at $q=q_a$ with
$B=M$, and inserting \eqref{mi32-block-weak}, gives
\[
 \bigl(\E\norm{Y[V_a,V_a]}^{2q_a}\bigr)^{1/(2q_a)}
 \le e\cdot2B_\alpha\bigl\{M+6\alpha^4J_\alpha^{10}D\bigr\}=L,
\]
by \eqref{mi32-L-constant}. Finally $q_a\ge2\pi_{a+1}
=2\cdot5^{a+1}\log2\ge5^{a+1}\ge a+2$, since $\log2>\tfrac12$.
\end{proof}

\begin{lemma}[A maximum of growing moments]
\label{mi32-max-lemma}
Let $Z_1,\ldots,Z_m\ge0$ and let $p_i\ge i+2$ be integers with
$\norm{Z_i}_{L_{p_i}}\le L$ for a constant $L>0$. Then
$\E\max_iZ_i\le3L$, with no dependence on $m$ and no independence
assumption.
\end{lemma}
\begin{proof}
For $z\ge0$, $t>0$ and $p\ge1$ one has
$z\le t+z^p/t^{p-1}$: this is clear if $z\le t$, and otherwise
$z\cdot t^{p-1}\le z\cdot z^{p-1}$. With $t=2L$,
\[
 \max_iZ_i\le2L+\sum_{i=1}^m\frac{Z_i^{p_i}}{(2L)^{p_i-1}} .
\]
Taking expectations and using $\E Z_i^{p_i}\le L^{p_i}$ and
$L^{p_i}/(2L)^{p_i-1}=L\,2^{-(p_i-1)}\le L\,2^{-(i+1)}$ gives
$\E\max_iZ_i\le2L+L\sum_{i\ge1}2^{-(i+1)}\le3L$.
\end{proof}

\subsection{Proof of Theorem~\ref{mi32-symmetric-deletion}}

If $L=0$ in \eqref{mi32-block-bound-constant}, replace $L$ by
$L+\epsilon$ throughout and let $\epsilon\downarrow0$ at the end; all
the bounds below are continuous in $L$. So assume $L>0$.

The mask $Y^{\rm nL}$ is supported on $S_1\times S_1$ with
$|S_1|\le b_1=31$ by \eqref{mi32-nested-card}. Hence
\begin{equation}\label{mi32-nearlow-bound}
 \E\norm{Y^{\rm nL}}
 \le\bigl(\E\norm{Y^{\rm nL}}_F^2\bigr)^{1/2}
 =\Bigl(\sum_{i\in S_1}\sum_{j\in S_1}v_{ij}\Bigr)^{1/2}
 \le\bigl(|S_1|M^2\bigr)^{1/2}\le\sqrt{31}\,M .
\end{equation}
This is the only block for which no deletion budget is available, and it
is also the only one whose dimension is bounded by an absolute
constant.

For $Y^{\rm nH}$ apply Lemma~\ref{mi32-block-norm-lemma} with
$c(i)=\lfloor\ell_i/2\rfloor$ and
$Z_m=\norm{Y^{\rm nH}[c^{-1}(m),c^{-1}(m)]}$, so that
$\norm{Y^{\rm nH}}=\max_mZ_m$. The block at $m=0$ vanishes, and for
$m\ge1$ it equals $Y[V_{2m},V_{2m}]$, to which
Proposition~\ref{mi32-block-proposition} applies at $a=2m\ge1$ with
$p_m=2q_{2m}\ge2m+2\ge m+2$. Lemma~\ref{mi32-max-lemma} therefore gives
$\E\norm{Y^{\rm nH}}\le3L$.

For $Y^{\rm nC}$ apply Lemma~\ref{mi32-block-norm-lemma} with
$c(i)=\lfloor(\ell_i+1)/2\rfloor$ and
$\sigma(i)=\ell_i\bmod2$. The block at $m=0$ vanishes, since there
$\ell_i=\ell_j=0$. For $m\ge1$ the two levels present in $c^{-1}(m)$
are $2m-1$ and $2m$, which have opposite parities, so the nearCross
condition $\ell_i\ne\ell_j$ is exactly $\sigma(i)\ne\sigma(j)$; the
second part of Lemma~\ref{mi32-block-norm-lemma} bounds that block by
$\norm{Y[V_{2m-1},V_{2m-1}]}$, which
Proposition~\ref{mi32-block-proposition} controls at $a=2m-1\ge1$ with
$p_m=2q_{2m-1}\ge2m+1\ge m+2$. Again
$\E\norm{Y^{\rm nC}}\le3L$.

Adding these to Proposition~\ref{mi32-far-proposition} through
\eqref{mi32-five-masks},
\begin{equation}\label{mi32-assembly}
 \E\norm{Y}\le\sqrt{31}\,M+6L+2M\sqrt{1+\sqrt2\,\alpha^2}.
\end{equation}
By \eqref{mi32-block-bound-constant} and \eqref{mi32-L-constant},
\[
 6L=24e\cdot5^{1/4}\sqrt3\,\alpha^2M
      +144e\cdot5^{1/4}\sqrt3\,\alpha^6J_\alpha^{10}D,
\]
so the right-hand side of \eqref{mi32-assembly} is
$c_1M+c_2D$ with $c_1,c_2\ge0$ summing to $C^{\rm sym}_\alpha$ as in
\eqref{mi32-Csym}. Since $M,D\ge0$, this is at most
$C^{\rm sym}_\alpha(M+D)$. \qed

\subsection{Symmetrization without changing the small exponent}
\label{mi32-symmetrization}

For general independent centered entries $X$, let $X'$ be an
independent copy and put $Y=X-X'$. For every $p\ge1$, conditional
Jensen and the triangle inequality give
\[
 \norm{X_{ij}}_{L_p}\le\norm{Y_{ij}}_{L_p},\qquad
 \norm{Y_{ij}}_{L_{2p}}\le2\alpha\norm{Y_{ij}}_{L_p}.
\]
Thus $Y$ has independent entries that are symmetric in law and satisfy
\eqref{mi32-regularity} with $2\alpha$; moreover
$\E Y_{ij}^2=2\E X_{ij}^2$, so $M(Y)=\sqrt2M(X)$, and
$\E\norm{X}\le\E\norm{Y}$ by Jensen applied to the convex function
$\norm{\cdot}$ in the copy coordinate.

For each fixed $I,s,t$, write the bilinear sum as $S$ and its
independent copy as $S'$. If $p\ge1$ then
$\norm{S-S'}_{L_p}\le2\norm{S}_{L_p}$. If $\log2\le p<1$,
subadditivity of $x\mapsto x^p$ gives instead
\[
 \norm{S-S'}_{L_p}\le2^{1/p}\norm{S}_{L_p}\le e\norm{S}_{L_p},
\]
since $2^{1/p}\le2^{1/\log2}=e$. Because $2\le e$, the factor $e$ covers
both ranges. Taking the same bilinear suprema, then the same
deterministic minimum over each budget --- the minimizer for $X$ is a
feasible set for $Y$ --- and then the budget maximum gives
$D(Y)\le eD(X)$. Every budget keeps its exact order $\log(k+1)$, and the
same original set is deleted from both axes throughout.

Applying Theorem~\ref{mi32-symmetric-deletion} to $Y$ with parameter
$2\alpha$ and using $\sqrt2\le e$ therefore proves
\begin{equation}\label{mi32-reduction-final-upper}
 \E\norm{X}\le C^{\rm sym}_{2\alpha}
     \bigl\{\sqrt2M(X)+eD(X)\bigr\}
 \le e\,C^{\rm sym}_{2\alpha}\{M(X)+D(X)\},
\end{equation}
which is the upper comparison in Theorem~\ref{mi32-main-theorem} with
\begin{equation}\label{mi32-final-constant}
 C_\alpha=e\,C^{\rm sym}_{2\alpha}.
\end{equation}
\section{The known lower bound on the same matrix}
\label{mi32-lower-section}

For completeness we reproduce the argument of
\cite[Theorem~4.1 and Lemma~4.2]{LatalaSwiatkowski}, with its
small-order and independence details. In this section $X$ itself
has independent centered entries satisfying \eqref{mi32-regularity}.
There is no symmetrization or comparison to a second input matrix.
Nothing in this section is used for the upper bound, and none of it is
part of the formalized statement.

\subsection{The scalar growth input}

Unlike the upper bound, the lower bound genuinely uses a moment-growth
estimate from the literature. By
\cite[Corollary~1.4]{LatalaStrzelecka}, applied to the singleton family
of test vectors, there is $K_\alpha\ge1$ such that every deterministic
linear combination $S$ of independent centered variables satisfying
\eqref{mi32-regularity} obeys
\begin{equation}\label{mi32-lower-scalar-growth}
 \norm{S}_{L_p}\le K_\alpha(p/r)^\beta\norm{S}_{L_r},
 \qquad p\ge r\ge2,
 \qquad \beta=\max\{1/2,\log_2\alpha\},
\end{equation}
and the exponent $\beta$ is optimal
\cite[Remark~1.5]{LatalaStrzelecka}. The same estimate appears as
\cite[(27)]{LatalaSwiatkowski}. Note that
Proposition~\ref{mi32-real-doubling} is \emph{not} a substitute here:
it compares $2p$ with $p$, whereas the sparsification below needs the
growth rate in $p$ itself.

\subsection{Row and column variances}

Fix a row and put $Q=\sum_jX_{ij}^2$, $v=\E Q$.
Regularity at order two and Cauchy--Schwarz give
\[
 \E Q^2=\sum_{j,l}\E(X_{ij}^2X_{il}^2)
 \le\alpha^4\sum_{j,l}\E X_{ij}^2\E X_{il}^2
 =\alpha^4v^2.
\]
Paley--Zygmund implies
$\E\sqrt Q\ge(4\sqrt2\alpha^4)^{-1}\sqrt v$.
Since every row and column Euclidean norm is at most $\norm{X}$,
\begin{equation}\label{mi32-lower-variance}
 M(X)\le8\sqrt2\alpha^4\E\norm{X}.
\end{equation}
Rows or columns with zero variance cause no exception.

\subsection{A large weak moment has a small original support}

For $p\ge2$, let $R_X^{\rm sp}(p)$ be the bilinear supremum
restricted to vectors whose supports each have cardinality at most
$p^{2\beta}$. Then
\begin{equation}\label{mi32-lower-sparsification}
 R_X(p)\le R_X^{\rm sp}(p)+C_\alpha M(X).
\end{equation}
Here is a direct proof. For a unit vector $z$, set
$I_p(z)=\{i:|z_i|\ge p^{-\beta}\}$, so that
$|I_p(z)|\le p^{2\beta}$. Split a bilinear sum into the part on
$I_p(s)\times I_p(t)$, the part with $j\notin I_p(t)$, and the
remaining part with $i\notin I_p(s)$ and $j\in I_p(t)$.
The first part is controlled by $R_X^{\rm sp}(p)$. For the second,
independence and zero means give
\[
 \left\|\sum_{i,\,j\notin I_p(t)}s_it_jX_{ij}\right\|_2^2
 \le p^{-2\beta}\sum_i s_i^2\sum_j\E X_{ij}^2
 \le p^{-2\beta}\sigma_r(X)^2.
\]
Apply \eqref{mi32-lower-scalar-growth} with lower order two
to bound its $L_p$ moment by $C_\alpha\sigma_r(X)$.
The third part is bounded in the same way by
$C_\alpha\sigma_c(X)$. The triangle inequality proves
\eqref{mi32-lower-sparsification}, also on every fixed deleted
submatrix with the same constant. The union of the two sparse
vector supports uses at most $2p^{2\beta}$ original indices.

\subsection{A uniform scalar lower tail}

Choose $D_\alpha\ge e$ so large that
$\norm{S}_{L_{2r}}\le D_\alpha\norm{S}_{L_r}$ for $r\ge2$,
using \eqref{mi32-lower-scalar-growth}, and set
$\delta=(4\log D_\alpha)^{-1}\le1/4$.
For $P\ge2/\delta$, take $r=\delta P\ge2$.
Paley--Zygmund applied to $|S|^r$ and scalar moment comparison give
\[
 \begin{split}
 \Prob\{|S|\ge\tfrac12\norm{S}_{L_r}\}
 &\ge(1-2^{-r})^2
       \left(\frac{\norm{S}_{L_r}}{\norm{S}_{L_{2r}}}\right)^{2r}
       \ge\tfrac14 e^{-P/2},\\
 \norm{S}_{L_r}&\ge K_\alpha^{-1}\delta^\beta\norm{S}_{L_P}.
 \end{split}
\]
Thus, with $c_0=\delta^\beta/(2K_\alpha)>0$,
\begin{equation}\label{mi32-lower-tail}
 \Prob\{|S|\ge c_0\norm{S}_{L_P}\}\ge\tfrac14e^{-P/2}.
\end{equation}
Only the scalar comparison of the original independent entries is used.

\subsection{Disjoint witnesses for each deletion budget}

Fix $1\le k\le n$ and write
\[
 P=\log(k+1),\qquad
 \gamma_k=\min_{|I|\le k}R_{X,I}(P).
\]
We prove $\gamma_k\le C_\alpha\E\norm{X}$ uniformly in $k$.
Choose $k_0(\alpha)$ large enough that, for $k\ge k_0$,
\[
 P\ge2/\delta,\qquad 2P^{2\beta}\le\sqrt k,\qquad k\ge4.
\]
For $k<k_0$, moment monotonicity and
\eqref{mi32-lower-scalar-growth} show
\[
 \gamma_k\le R_X(P)
 \le C_\alpha\max_{ij}\norm{X_{ij}}_2
 \le C_\alpha M(X).
\]
For $P<2$, only monotonicity to order two is used, so this covers
$k=1$ with its unchanged exponent $\log2$. Equation
\eqref{mi32-lower-variance} completes these finitely many orders.

Now let $k\ge k_0$. If $\gamma_k\le2C_\alpha M(X)$ with the
constant from \eqref{mi32-lower-sparsification}, the variance
bound again suffices. Otherwise, for every deterministic
$|I|\le k$, sparsification in $X[I^c,I^c]$ supplies unit vectors
$s,t$ supported outside $I$ on a common set $J$ of at most
$2P^{2\beta}\le\sqrt k$ indices, with
\[
 \left\|\sum_{i,j\in J}s_it_jX_{ij}\right\|_{L_P}
       \ge\gamma_k/2.
\]
The suprema are attained: there are finitely many possible supports,
and the moment functional is continuous on each compact product of
unit balls. Alternatively an arbitrarily small slack gives the
same constants.

Apply this observation successively, deleting the union of the
previously selected sets. It constructs
$m=\lfloor\sqrt k\rfloor$ disjoint deterministic sets $J_1,\ldots,J_m$
and bilinear sums $S_1,\ldots,S_m$ with
$\norm{S_l}_{L_P}\ge\gamma_k/2$. At every stage the already
deleted union has at most $(m-1)\sqrt k\le k$ elements.
The random variables $S_l$ are independent because they use disjoint
collections of original entries. Each satisfies $|S_l|\le\norm{X}$
pointwise, and \eqref{mi32-lower-tail} gives
\[
 \Prob\{|S_l|\ge c_0\gamma_k/2\}
       \ge\frac1{4\sqrt{k+1}}.
\]
Consequently
\[
 \begin{split}
 \E\norm{X}
 &\ge\frac{c_0\gamma_k}{2}
   \left[1-\left(1-\frac1{4\sqrt{k+1}}\right)^m\right]\\
 &\ge\frac{c_0\gamma_k}{2}
       \left(1-e^{-1/(8\sqrt2)}\right).
 \end{split}
\]
In the last step we used $m\ge\sqrt k/2$ and $k\ge4$.
Taking the maximum over the original budgets proves
\begin{equation}\label{mi32-lower-deletion}
 D(X)\le C_\alpha\E\norm{X}.
\end{equation}
Combining \eqref{mi32-lower-variance} and
\eqref{mi32-lower-deletion} gives the known same-matrix lower
estimate. Together with \eqref{mi32-reduction-final-upper}, it
proves Theorem~\ref{mi32-main-theorem}.
For $n=1$, $D(X)=0$ and the direct comparison
$\norm{X_{11}}_2\le\alpha\E|X_{11}|$ also verifies the lower
endpoint. All deletion minima throughout this proof remove the
same deterministic original set from both axes.


\section{Scope and the information retained by the proof}
\label{mi32-scope}

The estimate concerns the expected norm of the same original matrix
throughout. Its constants depend only on the entry regularity parameter.
The deletion sets are deterministic for the law, use the same original
indices on both axes, and retain the exact moment order $\log(k+1)$,
including $\log2<1$ at $k=1$.

Three features of the local argument prevent unsupported reductions.
First, the magnitudes and their repeated powers are never replaced by
static entry heights or fresh copies. Second, sector orthogonality
comes from the original signs and current vertices, not from magnitude
monomials on disjoint rows. Third, the multiplier estimate is used only
on a cone preserved by the actual walk prefixes; it does not bound all
signed linear combinations in a polynomial span.

What is \emph{not} claimed. The constants are explicit but
deliberately conservative, and no attempt is made to optimize them; in
particular Remark~\ref{mi32-sharp-cone-remark} already identifies one
available improvement that we do not take. No efficient algorithm for
evaluating $D(X)$ or $R_X(p)$ is claimed. The moment estimates supply
upper-tail bounds through Markov's inequality, but do not by themselves
identify matching quantiles, and matching moments are not used to claim
matching quantiles. Arbitrary noncommuting coefficient insertions can
destroy the nonnegative monomial cone, so applications with such
insertions need their own preserved-coefficient argument. Finally, the
lower comparison of Section~\ref{mi32-lower-section} is quoted from
\cite{LatalaSwiatkowski} and is not proved here in any new way.

\section{The formal verification}
\label{mi32-formalization}

The upper comparison has been formalized in Lean~4
\cite{LeanFour} on top of \texttt{mathlib} \cite{Mathlib}. This section
records what that does and does not establish.

\subsection{The formal statement}

The formalization fixes the quantities of
Section~\ref{mi32-introduction} as ordinary definitions over an
arbitrary probability space: the moment functional
$\norm{\cdot}_{L_p}$ as a real power of a Bochner integral, meaningful
for $p<1$; the operator norm as the norm of the Euclidean linear map;
$M$ and $D$ exactly as in \eqref{mi32-variances} and \eqref{mi32-D},
with one deterministic set deleted from both axes and the order
$\log(k+1)$ unchanged. The compared statement is
\begin{quote}\ttfamily\small
theorem main\_upper ($\alpha$ : $\mathbb R$) (h$\alpha$ : 1 $\le$ $\alpha$) :\\
\phantom{xx}$\exists$ C : $\mathbb R$, 0 < C $\wedge$ UpperBoundAt $\alpha$ C
\end{quote}
where \texttt{UpperBoundAt} quantifies over every dimension $n\ge1$,
every measurable space, every probability measure on it, and every
matrix of independent centered entries with all absolute moments finite
satisfying \eqref{mi32-regularity} at every real order $r\ge1$. No
finite-law substitute, conditional comparison, or assumed contraction
occurs anywhere, and no estimate is introduced as a hypothesis or as a
custom axiom. The transitive axiom audit of the target reports exactly
\texttt{propext}, \texttt{Classical.choice} and \texttt{Quot.sound}.

A separate check confirms that the hypothesis class is not vacuous: an
independent fair-sign matrix satisfies the hypotheses at the extreme
parameter $\alpha=1$ and has $M=\sqrt n>0$, so
\eqref{mi32-main-bound} is a genuine constraint rather than a statement
about an empty class.

\subsection{Correspondence with this manuscript}

\begin{center}\small
\begin{tabular}{@{}ll@{}}
\toprule
Statement here & Formal name in the \texttt{MI32} namespace\\
\midrule
Theorem~\ref{mi32-main-theorem}, upper half & \texttt{main\_upper}\\
Lemma~\ref{mi32-hilbert-boolean} & \texttt{hilbert\_walsh\_fourth\_le}\\
Lemma~\ref{mi32-product-moment-lemma} & \texttt{PositivePolynomial.tensorMoment\_add\_le}\\
Lemma~\ref{mi32-cone-lemma} & \texttt{SymmetricPositiveCone.fourth\_norm\_moment}\\
 & \texttt{SymmetricConeInterpolation.moment\_holder\_order\_le}\\
Corollary~\ref{mi32-scalar-doubling} & \texttt{PositiveLinearPowers.signed\_linear\_even\_root\_doubling}\\
Lemma~\ref{mi32-even-comparison} & \texttt{SymmetricMomentComparison.sum\_even\_moment\_le}\\
Theorem~\ref{mi32-hilbert-comparison} & \texttt{HilbertComparison.independent\_coordinate\_moment\_le}\\
Proposition~\ref{mi32-real-doubling} & \texttt{SymmetricLinearAllOrders.linear\_root\_doubling}\\
Lemma~\ref{mi32-thin-lemma} & \texttt{ThinMatrix.transposeOperator\_moment\_four\_mul\_le}\\
Proposition~\ref{mi32-contraction-proposition} & \texttt{RawCreatorContraction.creator\_l2\_le}\\
 & \texttt{RawAnnihilatorContraction.annihilator\_l2\_le}\\
Theorem~\ref{mi32-local-theorem} & \texttt{LocalSymmetricMoment.operator\_moment\_le}\\
Corollary~\ref{mi32-local-mean} & \texttt{LocalLogMoment.mean\_spectralNorm\_le\_literal}\\
Lemma~\ref{mi32-schedule-lemma} & \texttt{DeletionScheduleArithmetic}\\
Lemma~\ref{mi32-nested-lemma} & \texttt{NestedDeletionFamily}\\
Lemma~\ref{mi32-partition-lemma} & \texttt{DeletionMasks.entry\_sum}\\
Proposition~\ref{mi32-far-proposition} & \texttt{FarMaskMean.mean\_farLow\_le}, \texttt{mean\_farUp\_le}\\
Lemma~\ref{mi32-block-norm-lemma} & \texttt{FiberBlockNorm.norm\_le\_pi\_blockSub}\\
 & \texttt{FiberBlockNorm.norm\_offDiagonal\_le}\\
Proposition~\ref{mi32-block-proposition} & \texttt{NearBlockMoment.block\_moment\_le}\\
Lemma~\ref{mi32-max-lemma} & \texttt{BlockMomentMaximum.mean\_max\_le}\\
Theorem~\ref{mi32-symmetric-deletion} & \texttt{SymmetricDeletion.symmetricUpperBoundAt}\\
Section~\ref{mi32-symmetrization} & \texttt{SymmetrizationReduction.upperBoundAt\_of\_symmetric}\\
\bottomrule
\end{tabular}
\end{center}

\subsection{Differences between the two proofs, and limits of the check}

The formal proof follows this manuscript, with two deliberate
differences, both already flagged where they occur.

\begin{enumerate}
\item The cone constant. The formalization retains an earlier,
conservative form of the product-moment step, so the certified cone
comparison carries $H_\alpha=\sqrt3\,\alpha^2$ rather than the
$\sqrt3\,\alpha$ that Lemma~\ref{mi32-product-moment-lemma} in fact
gives (Remark~\ref{mi32-sharp-cone-remark}). We have quoted the
certified value everywhere, so that every displayed constant in this
manuscript is machine-checked; the sharp value would improve each of
them.
\item The dimension factor in Theorem~\ref{mi32-local-theorem} counts
both sides of the bipartite dilation, giving $2n$ where the trace
comparison between the two diagonal blocks gives $n$. At the orders
used this costs at most $\sqrt2$.
\end{enumerate}

The verification is a check of the deduction, not of the modelling.
It certifies that the Lean statement follows from the axioms of
Lean's type theory; whether that statement faithfully expresses
problem MI-32 is a question about the definitions, which a reader must
judge from the Challenge file and from
Section~\ref{mi32-introduction}. The lower comparison of
Section~\ref{mi32-lower-section} is not formalized, so
Theorem~\ref{mi32-main-theorem} is machine-checked only in its upper
half. Independent re-verification by the Comparator and NanoDa tools
of the Palomar protocol, an editorial review, and registration are
separate steps. All three have now been carried out: commit
	exttt{762bd5e} of the development repository is registered as Palomar
entry 	exttt{PALOMAR-2026-09-14-000006}, version~1, whose mechanical
report records Comparator, NanoDa and a sandboxed kernel replay
admitting only the three standard axioms, and whose editorial review
returned no blocking problem. That review is automated; it is not
external human peer review, and registration asserts neither novelty
nor endorsement.
% MI32 BODY END

\begin{thebibliography}{99}\raggedright

\bibitem{BandeiraVanHandel}
A.~S.~Bandeira and R.~van Handel.
\emph{Sharp nonasymptotic bounds on the norm of random matrices with
independent entries}. Ann.\ Probab.\ \textbf{44} (2016), 2479--2506.
\url{https://arxiv.org/abs/1408.6185}.

\bibitem{Latala2005}
R.~Lata\l a.
\emph{Some estimates of norms of random matrices}.
Proc.\ Amer.\ Math.\ Soc.\ \textbf{133} (2005), 1273--1282.

\bibitem{Latala2024}
R.~Lata\l a.
\emph{On the spectral norm of Rademacher matrices}.
Preprint, 2024. \url{https://arxiv.org/abs/2405.13656}.

\bibitem{LatalaStrzelecka}
R.~Lata\l a and M.~Strzelecka.
\emph{Comparison of weak and strong moments for vectors with independent
coordinates}. Mathematika \textbf{64} (2018), 211--229.
\url{https://arxiv.org/abs/1612.02407}.

\bibitem{LatalaSwiatkowski}
R.~Lata\l a and W.~\'Swi\k{a}tkowski.
\emph{Norms of randomized circulant matrices}.
Electron.\ J.\ Probab.\ \textbf{27} (2022), paper~80, 1--23.
Theorem, lemma and remark numbers refer to arXiv version~2.
\url{https://arxiv.org/abs/2106.03139v2}.

\bibitem{LatalaVanHandelYoussef}
R.~Lata\l a, R.~van Handel, and P.~Youssef.
\emph{The dimension-free structure of nonhomogeneous random matrices}.
Invent.\ Math.\ \textbf{214} (2018), 1031--1080.
\url{https://arxiv.org/abs/1711.00807}.

\bibitem{LeanFour}
L.~de Moura and S.~Ullrich.
\emph{The Lean 4 theorem prover and programming language}.
In: Automated Deduction --- CADE~28, Lecture Notes in Computer Science
\textbf{12699}, Springer, 2021, 625--635.

\bibitem{Mathlib}
The \texttt{mathlib} Community.
\emph{The Lean mathematical library}.
In: Proceedings of the 9th ACM SIGPLAN International Conference on
Certified Programs and Proofs (CPP~2020), 367--381.

\bibitem{Meller}
R.~Meller.
\emph{Spectral norm of matrices with independent entries up to
polyloglog}. Preprint, 2025.
\url{https://arxiv.org/abs/2512.23673}.

\bibitem{MI32}
\emph{MI-32: Spectral norms of independent entries with regular moment
growth}. Open Problems in Numerical Linear Algebra.
\url{https://github.com/ajt60gaibb/OpenProblemsInNLA/blob/main/matrix-inequalities-and-norms/MI-32/README.md}.

\bibitem{ODonnellBoolean}
R.~O'Donnell.
\emph{Analysis of Boolean Functions}.
Cambridge University Press, 2014.

\bibitem{RiemerSchutt}
S.~Riemer and C.~Sch\"utt.
\emph{On the expectation of the norm of random matrices with
non-identically distributed entries}.
Electron.\ J.\ Probab.\ \textbf{18} (2013), paper~29, 1--13.

\bibitem{Seginer}
Y.~Seginer.
\emph{The expected norm of random matrices}.
Combin.\ Probab.\ Comput.\ \textbf{9} (2000), 149--166.

\bibitem{VanHandelSurvey}
R.~van Handel.
\emph{Structured random matrices}.
In: Convexity and Concentration, IMA Volumes in Mathematics and its
Applications \textbf{161}, Springer, 2017, 107--156.

\end{thebibliography}
\end{document}
